%
% $Header: /home/cvs/root/haskell-report/report/exps.verb,v 1.20 2003/01/13 13:08:55 simonpj Exp $
%
%*section 3
%**<title>The Haskell 98 Report: Expressions</title>
%**~header
\chapter{Expressions}\index{expression}
\label{expressions}

In this chapter, we describe the syntax and informal semantics of
\Haskell{} {\em expressions}, including their translations into the
\Haskell{} kernel, where appropriate.  Except in the case of \mbox{\tt let}
expressions, these translations preserve both the static and dynamic
semantics.  Free variables and constructors used in these translations
always refer to entities defined by the \mbox{\tt Prelude}.  For example,
``\mbox{\tt concatMap}'' used in the translation of list comprehensions
(Section~\ref{list-comprehensions}) means the \mbox{\tt concatMap} defined by
the \mbox{\tt Prelude}, regardless of whether or not the identifier ``\mbox{\tt concatMap}'' is in
scope where the list comprehension is used, and (if it is in scope)
what it is bound to.

\begin{flushleft}\it\begin{tabbing}
\hspace{0.6in}\=\hspace{3.1in}\=\kill
$\it exp$\>\makebox[3.5em]{$\rightarrow$}$\it \hprime{infixexp\ \makebox{\tt ::}\ [context\ \makebox{\tt =>}]\ type}$\>\makebox[3em]{}$\it (\tr{expression\ type\ signature})$\\ 
$\it $\>\makebox[3.5em]{$|$}$\it \hprime{infixexp}$\\ 
$\it $\\ 
$\it \hprime{infixexp}$\>\makebox[3.5em]{$\rightarrow$}$\it \hprime{lexp\ qop\ infixexp}$\>\makebox[3em]{}$\it (\tr{infix\ operator\ application})$\\ 
$\it $\>\makebox[3.5em]{$|$}$\it \hprime{\makebox{\tt -}\ infixexp}$\>\makebox[3em]{}$\it (\tr{prefix\ negation})$\\ 
$\it $\>\makebox[3.5em]{$|$}$\it \hprime{lexp}$\\ 
$\it $\\ 
$\it \hprime{lexp}$\>\makebox[3.5em]{$\rightarrow$}$\it \makebox{\tt {\char'134}}\ apat_1\ \ldots \ apat_n\ \makebox{\tt ->}\ exp$\>\makebox[3em]{}$\it (\tr{lambda\ abstraction},\ n\geq 1)$\\ 
$\it $\>\makebox[3.5em]{$|$}$\it \makebox{\tt let}\ decls\ \makebox{\tt in}\ exp$\>\makebox[3em]{}$\it ({\tr{let\ expression}})$\\ 
$\it $\>\makebox[3.5em]{$|$}$\it \makebox{\tt if}\ exp\ \hprime{[\makebox{\tt ;}]}\ \makebox{\tt then}\ exp\ \hprime{[\makebox{\tt ;}]}\ \makebox{\tt else}\ exp$\>\makebox[3em]{}$\it (\tr{conditional})$\\ 
$\it $\>\makebox[3.5em]{$|$}$\it \makebox{\tt case}\ exp\ \makebox{\tt of}\ \makebox{\tt {\char'173}}\ alts\ \makebox{\tt {\char'175}}$\>\makebox[3em]{}$\it (\tr{case\ expression})$\\ 
$\it $\>\makebox[3.5em]{$|$}$\it \makebox{\tt do}\ \makebox{\tt {\char'173}}\ stmts\ \makebox{\tt {\char'175}}$\>\makebox[3em]{}$\it (\tr{do\ expression})$\\ 
$\it $\>\makebox[3.5em]{$|$}$\it fexp$\\ 
$\it fexp$\>\makebox[3.5em]{$\rightarrow$}$\it [fexp]\ aexp$\>\makebox[3em]{}$\it (\tr{function\ application})$\\ 
$\it $\\ 
$\it aexp$\>\makebox[3.5em]{$\rightarrow$}$\it qvar$\>\makebox[3em]{}$\it (\tr{variable})$\\ 
$\it $\>\makebox[3.5em]{$|$}$\it gcon$\>\makebox[3em]{}$\it (\tr{general\ constructor})$\\ 
$\it $\>\makebox[3.5em]{$|$}$\it literal$\\ 
$\it $\>\makebox[3.5em]{$|$}$\it \makebox{\tt (}\ exp\ \makebox{\tt )}$\>\makebox[3em]{}$\it (\tr{parenthesized\ expression})$\\ 
$\it $\>\makebox[3.5em]{$|$}$\it \makebox{\tt (}\ exp_1\ \makebox{\tt ,}\ \ldots \ \makebox{\tt ,}\ exp_k\ \makebox{\tt )}$\>\makebox[3em]{}$\it (\tr{tuple},\ k\geq 2)$\\ 
$\it $\>\makebox[3.5em]{$|$}$\it \makebox{\tt [}\ exp_1\ \makebox{\tt ,}\ \ldots \ \makebox{\tt ,}\ exp_k\ \makebox{\tt ]}$\>\makebox[3em]{}$\it (\tr{list},\ k\geq 1)$\\ 
$\it $\>\makebox[3.5em]{$|$}$\it \makebox{\tt [}\ exp_1\ [\makebox{\tt ,}\ exp_2]\ \makebox{\tt ..}\ [exp_3]\ \makebox{\tt ]}$\>\makebox[3em]{}$\it (\tr{arithmetic\ sequence})$\\ 
$\it $\>\makebox[3.5em]{$|$}$\it \makebox{\tt [}\ exp\ \makebox{\tt |}\ qual_1\ \makebox{\tt ,}\ \ldots \ \makebox{\tt ,}\ qual_n\ \makebox{\tt ]}$\>\makebox[3em]{}$\it (\tr{list\ comprehension},\ n\geq 1)$\\ 
$\it $\>\makebox[3.5em]{$|$}$\it \makebox{\tt (}\ \hprime{infixexp\ qop}\ \makebox{\tt )}$\>\makebox[3em]{}$\it (\tr{left\ section})$\\ 
$\it $\>\makebox[3.5em]{$|$}$\it \makebox{\tt (}\ \hprime{qop_{\langle\makebox{\tt -}\rangle}\ infixexp}\ \makebox{\tt )}$\>\makebox[3em]{}$\it (\tr{right\ section})$\\ 
$\it $\>\makebox[3.5em]{$|$}$\it qcon\ \makebox{\tt {\char'173}}\ fbind_1\ \makebox{\tt ,}\ \ldots \ \makebox{\tt ,}\ fbind_n\ \makebox{\tt {\char'175}}$\>\makebox[3em]{}$\it (\tr{labeled\ construction},\ n\geq 0)$\\ 
$\it $\>\makebox[3.5em]{$|$}$\it aexp_{\langle{}qcon\rangle{}}\ \makebox{\tt {\char'173}}\ fbind_1\ \makebox{\tt ,}\ \ldots \ \makebox{\tt ,}\ fbind_n\ \makebox{\tt {\char'175}}$\>\makebox[3em]{}$\it (\tr{labeled\ update},\ n\ \geq \ 1)$\\ 
$\it $
\end{tabbing}\end{flushleft}
\indexsyn{exp}%
\indexsyn{infixexp}%
\indexsyn{lexp}%
\indexsyn{aexp}%
\indexsyn{fexp}%

%       Removed Aug 2001: more misleading than helpful. SLPJ
% As an aid to understanding this grammar,
% Table~\ref{syntax-precedences} shows the relative precedence of
% expressions, patterns and definitions, plus an extended associativity.
% \mbox{$\it -$} indicates that the item is non-associative.
% 
% \begin{table}[tb]
% \[
% \centerline{
% \begin{tabular}{|l|c|}\hline
% Item                                          &       Associativity \\ 
% \hline
%                                               &                       \\
% simple terms, parenthesized terms               &     --              \\
% irrefutable patterns (\mbox{\tt {\char'176}})                      &       --            \\
% as-patterns ({\tt @})                          &       right         \\
% function application                          &       left            \\
% \mbox{\tt do}, \mbox{\tt if}, \mbox{\tt let}, lambda(\mbox{\tt {\char'134}}), \mbox{\tt case} (leftwards)&   right           \\
% \mbox{\tt case} (rightwards)                           &       right           \\
%                                               &                       \\
% infix operators, prec. 9                      &       as defined      \\
% \ldots                                                &       \ldots          \\
% infix operators, prec. 0                      &       as defined      \\
%                                               &                       \\
% function types (\mbox{\tt ->})                         &       right           \\
% contexts (\mbox{\tt =>})                                       &       --              \\
% type constraints (\mbox{\tt ::})                               &       --              \\
% \mbox{\tt do}, \mbox{\tt if}, \mbox{\tt let}, lambda(\mbox{\tt {\char'134}}) (rightwards)   &       right           \\
% sequences (\mbox{\tt ..})                              &       --              \\
% generators (\mbox{\tt <-})                             &       --              \\
% grouping (\mbox{\tt ,})                                        &       n-ary           \\
% guards (\mbox{\tt |})                                  &       --              \\
% case alternatives (\mbox{\tt ->})                      &       --              \\
% definitions (\mbox{\tt =})                             &       --              \\
% separation (\mbox{\tt ;})                              &       n-ary           \\ 
% \hline
% \end{tabular}
% }
% \]
% %**<div align=center> <h4>Table 1</h4> </div>
% \ecaption{Precedence of expressions, patterns, definitions (highest to lowest)}
% \label{syntax-precedences}
% \end{table}

Expressions involving infix operators are disambiguated by the
operator's fixity (see Section~\ref{fixity}).  Consecutive
unparenthesized operators with the same precedence must both be either
left or right associative to avoid a syntax error.
Given an unparenthesized expression ``\mbox{$\it x\ qop^{(a,i)}\ y\ qop^{(b,j)}\ z$}'' 
\hprime{(where \mbox{$\it qop^{(a,i)}$} means an operator with associativity \mbox{$\it a$} and
precedence \mbox{$\it i$})}, parentheses must be added around either ``\mbox{$\it x\\
\it qop^{(a,i)}\ y$}'' or ``\mbox{$\it y\ qop^{(b,j)}\ z$}'' when \mbox{$\it i=j$} unless \mbox{$\it a=b={\rm\\
\it l}$} or \mbox{$\it a=b={\rm\ r}$}.

An example algorithm for resolving expressions involving infix
operators is given in Section~\ref{fixity-resolution}.

Negation\index{negation} is the only prefix operator in
\Haskell{}; it has the same precedence as the infix \mbox{\tt -} operator
defined in the Prelude (see Section~\ref{fixity}, Figure~\ref{prelude-fixities}).

The grammar is ambiguous regarding the extent of lambda abstractions,
let expressions, and conditionals.  The ambiguity is resolved by the
meta-rule that each of these constructs extends as far to the right as
possible.  

%       I can't make head or tail of this para, so
%       I'm just deleting it.  SLPJ Dec 98
% The separation of function arrows from case alternatives solves
% the ambiguity that otherwise arises when an unparenthesized
% function type is used in an expression, such as the guard in a case
% expression.

Sample parses are shown below.
\[\bt{|l|l|}%%b
\hline
This                                & Parses as                             \\
\hline
\mbox{\tt f\ x\ +\ g\ y}                         & \mbox{\tt (f\ x)\ +\ (g\ y)}                       \\
\mbox{\tt -\ f\ x\ +\ y}                         & \mbox{\tt (-\ (f\ x))\ +\ y}                       \\
\mbox{\tt let\ {\char'173}\ ...\ {\char'175}\ in\ x\ +\ y}              & \mbox{\tt let\ {\char'173}\ ...\ {\char'175}\ in\ (x\ +\ y)}              \\
\mbox{\tt z\ +\ let\ {\char'173}\ ...\ {\char'175}\ in\ x\ +\ y}          & \mbox{\tt z\ +\ (let\ {\char'173}\ ...\ {\char'175}\ in\ (x\ +\ y))}        \\
\mbox{\tt f\ x\ y\ ::\ Int}                      & \mbox{\tt (f\ x\ y)\ ::\ Int}                      \\
\mbox{\tt {\char'134}\ x\ ->\ a+b\ ::\ Int}                 & \mbox{\tt {\char'134}\ x\ ->\ ((a+b)\ ::\ Int})               \\
\hline\et\]

For the sake of clarity, the rest of this section will assume that
expressions involving infix operators have been resolved according to
the fixities of the operators.

\section{Errors}
\label{basic-errors}\index{error}
Errors during expression evaluation, denoted by \mbox{$\it \bot$}\index{\mbox{$\it \bot$}} (``bottom''\index{bottom}),
are indistinguishable by a Haskell program from non-termination.  Since \Haskell{} is a
non-strict language, all \Haskell{} types include \mbox{$\it \bot$}.  That is, a value
of any type may be bound to a computation that, when demanded, results
in an error.  When evaluated, errors cause immediate program
termination and cannot be caught by the user.  The Prelude provides
two functions to directly 
cause such errors:
\bprog
{\small\begin{verbatim}
error     :: String -> a
undefined :: a
\end{verbatim}}
\eprog
\indextt{error}
\indextt{undefined}
A call to \mbox{\tt error} terminates execution of
the program and returns an appropriate error indication to the
operating system.  It should also display the string in some
system-dependent manner.  When \mbox{\tt undefined} is used, the error message
is created by the compiler.

Translations of \Haskell{} expressions use \mbox{\tt error} and \mbox{\tt undefined} to
explicitly indicate where execution time errors may occur.  The actual
program behavior when an error occurs is up to the implementation.
The messages passed to the \mbox{\tt error} function in these translations are
only suggestions; implementations may choose to display more or less
information when an error occurs.

\section{Variables, Constructors, Operators, and Literals}
\label{vars-and-lits}
%
\begin{flushleft}\it\begin{tabbing}
\hspace{0.6in}\=\hspace{3.1in}\=\kill
$\it aexp$\>\makebox[3.5em]{$\rightarrow$}$\it qvar$\>\makebox[3em]{}$\it (\tr{variable})$\\ 
$\it $\>\makebox[3.5em]{$|$}$\it gcon$\>\makebox[3em]{}$\it (\tr{general\ constructor})$\\ 
$\it $\>\makebox[3.5em]{$|$}$\it literal$
\end{tabbing}\end{flushleft}
\indexsyn{var}%
\indexsyn{con}%
\indexsyn{varop}%
\indexsyn{conop}%
\indexsyn{op}%
\begin{flushleft}\it\begin{tabbing}
\hspace{0.6in}\=\hspace{3.1in}\=\kill
$\it gcon$\>\makebox[3.5em]{$\rightarrow$}$\it \makebox{\tt ()}$\\ 
$\it $\>\makebox[3.5em]{$|$}$\it \makebox{\tt []}$\\ 
$\it $\>\makebox[3.5em]{$|$}$\it \makebox{\tt (,}\{\makebox{\tt ,}\}\makebox{\tt )}$\\ 
$\it $\>\makebox[3.5em]{$|$}$\it qcon$\\ 
$\it $\\ 
$\it var$\>\makebox[3.5em]{$\rightarrow$}$\it varid\ |\ \makebox{\tt (}\ varsym\ \makebox{\tt )}$\>\makebox[3em]{}$\it (\tr{variable})$\\ 
$\it qvar$\>\makebox[3.5em]{$\rightarrow$}$\it qvarid\ |\ \makebox{\tt (}\ qvarsym\ \makebox{\tt )}$\>\makebox[3em]{}$\it (\tr{qualified\ variable})$\\ 
$\it con$\>\makebox[3.5em]{$\rightarrow$}$\it conid\ |\ \makebox{\tt (}\ consym\ \makebox{\tt )}$\>\makebox[3em]{}$\it (\tr{constructor})$\\ 
$\it qcon$\>\makebox[3.5em]{$\rightarrow$}$\it qconid\ |\ \makebox{\tt (}\ gconsym\ \makebox{\tt )}$\>\makebox[3em]{}$\it (\tr{qualified\ constructor})$\\ 
$\it varop$\>\makebox[3.5em]{$\rightarrow$}$\it varsym\ |\ \bkqB\ varid\ \bkqA$\>\makebox[3em]{}$\it (\tr{variable\ operator})$\\ 
$\it qvarop$\>\makebox[3.5em]{$\rightarrow$}$\it qvarsym\ |\ \bkqB\ qvarid\ \bkqA$\>\makebox[3em]{}$\it (\tr{qualified\ variable\ operator})$\\ 
$\it conop$\>\makebox[3.5em]{$\rightarrow$}$\it consym\ |\ \bkqB\ conid\ \bkqA$\>\makebox[3em]{}$\it (\tr{constructor\ operator})$\\ 
$\it qconop$\>\makebox[3.5em]{$\rightarrow$}$\it gconsym\ |\ \bkqB\ qconid\ \bkqA$\>\makebox[3em]{}$\it (\tr{qualified\ constructor\ operator})$\\ 
$\it op$\>\makebox[3.5em]{$\rightarrow$}$\it varop\ |\ conop$\>\makebox[3em]{}$\it (\tr{operator})$\\ 
$\it qop$\>\makebox[3.5em]{$\rightarrow$}$\it qvarop\ |\ qconop$\>\makebox[3em]{}$\it (\tr{qualified\ operator})$\\ 
$\it gconsym$\>\makebox[3.5em]{$\rightarrow$}$\it \makebox{\tt :}\ |\ qconsym$
\end{tabbing}\end{flushleft}
\indexsyn{gcon}%
\indexsyn{var}%
\indexsyn{qvar}%
\indexsyn{con}%
\indexsyn{qcon}%
\indexsyn{varop}%
\indexsyn{qvarop}%
\indexsyn{conop}%
\indexsyn{qconop}%
\indexsyn{qop}%
\indexsyn{gconsym}%

\Haskell{} provides special syntax to support infix notation.
An {\em operator} is a function that can be applied using infix 
syntax (Section~\ref{operators}), or partially applied using a
{\em section} (Section~\ref{sections}).

An {\em operator} is either an {\em operator symbol}, such as \mbox{\tt +} or \mbox{\tt {\char'44}{\char'44}},
or is an ordinary identifier enclosed in grave accents (backquotes), such
as \bkqB\mbox{\tt op}\bkqA.  For example, instead of writing the prefix application
\mbox{\tt op\ x\ y}, one can write the infix application \mbox{\mbox{\tt x} \bkqB\mbox{\tt op}\bkqA\mbox{\tt \ y}}.
If no fixity\index{fixity}
declaration is given for \bkqB\mbox{\tt op}\bkqA{} then it defaults
to highest precedence and left associativity
(see Section~\ref{fixity}).

Dually, an operator symbol can be converted to an ordinary identifier
by enclosing it in parentheses.  For example, \mbox{\tt (+)\ x\ y} is equivalent
to \mbox{\tt x\ +\ y}, and \mbox{\tt foldr\ (*)\ 1\ xs} is equivalent to \mbox{\tt foldr\ ({\char'134}x\ y\ ->\ x*y)\ 1\ xs}.

%       This para is covered by Section 2.4 and 5.5.1
% A qualified name may only be used to refer to a variable or
% constructor imported from another module (see Section~\ref{import}), or
% defined at the top level, 
% but not in the definition of a new variable or constructor.  Thus
% \bprog
% let F.x = 1 in F.x   -- invalid
% \eprog
% incorrectly uses a qualifier in the definition of \mbox{\tt x}, regardless of
% the module containing this definition.  Qualification does not affect
% the nature of an operator: \mbox{\tt F.+} is an infix operator just as \mbox{\tt +} is.

Special syntax is used to name some constructors for some of the
built-in types, as found
in the production for \mbox{$\it gcon$} and \mbox{$\it literal$}.  These are described
in Section~\ref{basic-types}.

\index{number!translation of literals}
An integer literal represents the
application of the function \mbox{\tt fromInteger}\indextt{fromInteger} to the
appropriate value of type 
\mbox{\tt Integer}.  Similarly, a floating point literal stands for an application of
\mbox{\tt fromRational}\indextt{fromRational} to a value of type \mbox{\tt Rational} (that is, 
\mbox{\tt Ratio\ Integer}).

\outline{
\paragraph*{Translation:}
The integer literal \mbox{$\it i$} is equivalent to \mbox{\tt fromInteger} \mbox{$\it i$},
where \mbox{\tt fromInteger} is a method in class \mbox{\tt Num} (see Section
\ref{numeric-literals}).\indexclass{Num}


The floating point literal \mbox{$\it f$} is equivalent to \mbox{\tt fromRational}
(\mbox{$\it n$} \mbox{\tt Ratio.{\char'45}} \mbox{$\it d$}), where \mbox{\tt fromRational} is a method in class \mbox{\tt Fractional}
and \mbox{\tt Ratio.{\char'45}} constructs a rational from two integers, as defined in
the \mbox{\tt Ratio} library.\indexclass{Fractional}
The integers \mbox{$\it n$} and \mbox{$\it d$} are chosen so that \mbox{$\it n/d\ =\ f$}.
}


\section{Curried Applications and Lambda Abstractions}
\label{applications}
\label{lambda-abstractions}
\index{lambda abstraction}
\index{application}
%\index{function application|see{application}}
%
\begin{flushleft}\it\begin{tabbing}
\hspace{0.6in}\=\hspace{3.1in}\=\kill
$\it fexp$\>\makebox[3.5em]{$\rightarrow$}$\it [fexp]\ aexp$\>\makebox[3em]{}$\it (\tr{function\ application})$\\ 
$\it lexp$\>\makebox[3.5em]{$\rightarrow$}$\it \makebox{\tt {\char'134}}\ apat_1\ \ldots \ apat_n\ \makebox{\tt ->}\ exp$\>\makebox[3em]{}$\it (\tr{lambda\ abstraction},\ n\geq 1)$
\end{tabbing}\end{flushleft}
\indexsyn{exp}%
\indexsyn{fexp}%

\noindent
{\em Function application}\index{application} is written 
\mbox{$\it e_1\ e_2$}.  Application associates to the left, so the
parentheses may be omitted in \mbox{\tt (f\ x)\ y}.  Because \mbox{$\it e_1$} could
be a data constructor, partial applications of data constructors are
allowed. 

{\em Lambda abstractions} are written 
\mbox{$\it \makebox{\tt {\char'134}}\ p_1\ \ldots \ p_n\ \makebox{\tt ->}\ e$}, where the \mbox{$\it p_i$} are {\em patterns}.
An expression such as \mbox{\tt {\char'134}x:xs->x} is syntactically incorrect;
it may legally be written as \mbox{\tt {\char'134}(x:xs)->x}.

The set of patterns must be {\em linear}\index{linearity}%
\index{linear pattern}---no variable may appear more than once in the set.

\outline{\ifhtml{}{\small}
\paragraph*{Translation:}
The following identity holds:
\begin{center}
\bt{lcl}%
\struthack{17pt}%
\mbox{$\it \makebox{\tt {\char'134}}\ p_1\ \ldots \ p_n\ \makebox{\tt ->}\ e$}
         & \mbox{$\it =$} &
        \mbox{$\it \makebox{\tt {\char'134}}\ x_1\ \ldots \ x_n\ \makebox{\tt ->\ case\ (}x_1\makebox{\tt ,}\ \ldots \makebox{\tt ,}\ x_n\makebox{\tt )\ of\ (}p_1\makebox{\tt ,}\ \ldots \makebox{\tt ,}\ p_n\makebox{\tt )\ ->}\ e$}
\et
\end{center}
where the \mbox{$\it x_i$} are new identifiers.
}
Given this translation combined with the semantics of case
expressions and pattern matching described in
Section~\ref{case-semantics}, if the
pattern fails to match, then the result is \mbox{$\it \bot$}.

             
\section{Operator Applications}
\index{operator application}
%\index{operator application|hseealso{application}}
\label{operators}
%
\begin{flushleft}\it\begin{tabbing}
\hspace{0.6in}\=\hspace{3.1in}\=\kill
$\it \hprime{infixexp}$\>\makebox[3.5em]{$\rightarrow$}$\it \hprime{lexp\ qop\ infixexp}$\\ 
$\it $\>\makebox[3.5em]{$|$}$\it \makebox{\tt -}\ \hprime{infixexp}$\>\makebox[3em]{}$\it (\tr{prefix\ negation})$\\ 
$\it $\>\makebox[3.5em]{$|$}$\it \hprime{lexp}$\\ 
$\it qop$\>\makebox[3.5em]{$\rightarrow$}$\it qvarop\ |\ qconop$\>\makebox[3em]{}$\it (\tr{qualified\ operator})$
\end{tabbing}\end{flushleft}
\indexsyn{exp}%
\indexsyn{qop}%

\noindent
The form \mbox{$\it e_1\ qop\ e_2$} is the infix application of binary
operator\index{operator} \mbox{$\it qop$} to expressions \mbox{$\it e_1$} and \mbox{$\it e_2$}.  

The special
form \mbox{$\it \makebox{\tt -}e$} denotes prefix negation\index{negation}, the only
prefix operator in \Haskell{}, and is 
syntax for \mbox{$\it \makebox{\tt negate\ }(e)$}.\indextt{negate}  The binary \mbox{\tt -} operator
does not necessarily refer 
to the definition of \mbox{\tt -} in the Prelude; it may be rebound 
by the module system.  However, unary \mbox{\tt -} will always refer to the
\mbox{\tt negate} function defined in the Prelude.  There is no link between
the local meaning of the \mbox{\tt -} operator and unary negation.

Prefix negation has the same precedence as the infix operator \mbox{\tt -}
defined in the Prelude (see
Table~\ref{prelude-fixities}).  Because \mbox{\tt e1-e2} parses as an
infix application of the binary operator \mbox{\tt -}, one must write \mbox{\tt e1(-e2)} for
the alternative parsing.  Similarly, \mbox{\tt (-)} is syntax for 
\mbox{\tt ({\char'134}\ x\ y\ ->\ x-y)}, as with any infix operator, and does not denote 
\mbox{\tt ({\char'134}\ x\ ->\ -x)}---one must use \mbox{\tt negate} for that.

\outline{
\paragraph*{Translation:}
The following identities hold:
\begin{center}
\bt{lcl}%
\struthack{17pt}%
\mbox{$\it e_1\ op\ e_2$} & \mbox{$\it =$} & \mbox{$\it \makebox{\tt (}op\makebox{\tt )}\ e_1\ e_2$} \\
\mbox{$\it \makebox{\tt -}e$} & \mbox{$\it =$} & \mbox{$\it \makebox{\tt negate}\ (e)$}
\et
\end{center}
}

\section{Sections}
\index{section}
%\index{section|hseealso{operator application}}
\label{sections}
%
\begin{flushleft}\it\begin{tabbing}
\hspace{0.6in}\=\hspace{3.1in}\=\kill
$\it aexp$\>\makebox[3.5em]{$\rightarrow$}$\it \makebox{\tt (}\ \hprime{infixexp\ qop}\ \makebox{\tt )}$\>\makebox[3em]{}$\it (\tr{left\ section})$\\ 
$\it $\>\makebox[3.5em]{$|$}$\it \makebox{\tt (}\ \hprime{qop_{\langle\makebox{\tt -}\rangle}\ infixexp}\ \makebox{\tt )}$\>\makebox[3em]{}$\it (\tr{right\ section})$
\end{tabbing}\end{flushleft}
\indexsyn{infixexp}%
\indexsyn{qop}%

\noindent
{\em Sections} are written as \mbox{$\it \makebox{\tt (}\ op\ e\ \makebox{\tt )}$} or \mbox{$\it \makebox{\tt (}\ e\ op\ \makebox{\tt )}$}, where
\mbox{$\it op$} is a binary operator and \mbox{$\it e$} is an expression.  Sections are a
convenient syntax for partial application of binary operators.

Syntactic precedence rules apply to sections as follows.
\mbox{$\it \makebox{\tt (}op~e\makebox{\tt )}$} is legal if and only if \mbox{$\it \makebox{\tt (x}~op~e\makebox{\tt )}$} parses 
in the same way as \mbox{$\it \makebox{\tt (x}~op~\makebox{\tt (}e\makebox{\tt ))}$};
and similarly for  \mbox{$\it \makebox{\tt (}e~op\makebox{\tt )}$}.
For example, \mbox{\tt (*a+b)} is syntactically invalid, but \mbox{\tt (+a*b)} and
\mbox{\tt (*(a+b))} are valid.  Because \mbox{\tt (+)} is left associative, \mbox{\tt (a+b+)} is syntactically correct,
but \mbox{\tt (+a+b)} is not; the latter may legally be written as \mbox{\tt (+(a+b))}.
As another example, the expression
\bprog
{\small\begin{verbatim}
  (let n = 10 in n +)
\end{verbatim}}
\eprog
is invalid because, by the let/lambda meta-rule (Section~\ref{expressions}),
the expression
\bprog
{\small\begin{verbatim}
  (let n = 10 in n + x)
\end{verbatim}}
\eprog
parses as
\bprog
{\small\begin{verbatim}
  (let n = 10 in (n + x))
\end{verbatim}}
\eprog
rather than
\bprog
{\small\begin{verbatim}
  ((let n = 10 in n) + x)
\end{verbatim}}
\eprog
% This observation makes it easier to implement the let/lambda meta-rule
% (Section~\ref{expressions}) because once the operator has been seen it is clear that any
% legal parse must include the operator in the body of the \mbox{\tt let}.

Because \mbox{\tt -} is treated specially in the grammar,
\mbox{$\it \makebox{\tt (-}\ exp\makebox{\tt )}$} is not a section, but an application of prefix
negation,\index{negation} as
described in the preceding section.  However, there is a \mbox{\tt subtract}
function defined in the Prelude such that
\mbox{$\it \makebox{\tt (subtract}\ exp\makebox{\tt )}$} is equivalent to the disallowed section.
The expression \mbox{$\it \makebox{\tt (+\ (-}\ exp\makebox{\tt ))}$} can serve the same purpose.


% Changed to allow postfix operators.  That is, in (op x), we no
% longer add a \x -> which would require op to be binary instead
% of unary.

\outline{
\paragraph*{Translation:}
The following identities hold:
\begin{center}
\bt{lcl}%
\struthack{17pt}%
\mbox{$\it \makebox{\tt (}op\ e\makebox{\tt )}$} & \mbox{$\it =$} & \mbox{$\it \makebox{\tt {\char'134}}\ x\ \makebox{\tt ->}\ x\ op\ e$} \\
\mbox{$\it \makebox{\tt (}e\ op\makebox{\tt )}$} & \mbox{$\it =$} & \mbox{$\it \makebox{\tt {\char'134}}\ x\ \makebox{\tt ->}\ e\ op\ x$}
\et
\end{center}
where \mbox{$\it op$} is a binary operator, \mbox{$\it e$} is an expression, and \mbox{$\it x$} is a variable
that does not occur free in \mbox{$\it e$}.
}

\section{Conditionals}
\label{conditionals}\index{conditional expression}
%
\begin{flushleft}\it\begin{tabbing}
\hspace{0.6in}\=\hspace{3.1in}\=\kill
$\it lexp$\>\makebox[3.5em]{$\rightarrow$}$\it \makebox{\tt if}\ exp\ \hprime{[\makebox{\tt ;}]}\ \makebox{\tt then}\ exp\ \hprime{[\makebox{\tt ;}]}\ \makebox{\tt else}\ exp$
\end{tabbing}\end{flushleft}
\indexsyn{exp}%

%\indextt{if ... then ... else ...}
A {\em conditional expression}
\index{conditional expression}
has the form 
\mbox{$\it \makebox{\tt if}\ e_1\ \makebox{\tt then}\ e_2\ \makebox{\tt else}\ e_3$} and returns the value of \mbox{$\it e_2$} if the
value of \mbox{$\it e_1$} is \mbox{\tt True}, \mbox{$\it e_3$} if \mbox{$\it e_1$} is \mbox{\tt False}, and \mbox{$\it \bot$}
otherwise.

\outline{
\paragraph*{Translation:}
The following identity holds:
\begin{center}
\bt{lcl}%
\struthack{17pt}%
\mbox{$\it \makebox{\tt if}\ e_1\ \makebox{\tt then}\ e_2\ \makebox{\tt else}\ e_3$}  & \mbox{$\it =$} & \mbox{$\it \makebox{\tt case}\ e_1\ \makebox{\tt of\ {\char'173}\ True\ ->}\ e_2\ \makebox{\tt ;\ False\ ->}\ e_3\ \makebox{\tt {\char'175}}$}
\et
\end{center}
where \mbox{\tt True} and \mbox{\tt False} are the two nullary constructors from the 
type \mbox{\tt Bool}, as defined in the Prelude.  The type of \mbox{$\it e_1$} must be \mbox{\tt Bool};
\mbox{$\it e_2$} and \mbox{$\it e_3$} must have the same type, which is also the type of the
entire conditional expression.
}

\section{Lists}
\label{lists}
%
\begin{flushleft}\it\begin{tabbing}
\hspace{0.6in}\=\hspace{3.1in}\=\kill
$\it infixexp$\>\makebox[3.5em]{$\rightarrow$}$\it exp_1\ qop\ exp_2$\\ 
$\it aexp$\>\makebox[3.5em]{$\rightarrow$}$\it \makebox{\tt [}\ exp_1\ \makebox{\tt ,}\ \ldots \ \makebox{\tt ,}\ exp_k\ \makebox{\tt ]}$\>\makebox[3em]{}$\it (k\geq 1)$\\ 
$\it $\>\makebox[3.5em]{$|$}$\it gcon$\\ 
$\it gcon$\>\makebox[3.5em]{$\rightarrow$}$\it \makebox{\tt []}$\\ 
$\it $\>\makebox[3.5em]{$|$}$\it qcon$\\ 
$\it qcon$\>\makebox[3.5em]{$\rightarrow$}$\it \makebox{\tt (}\ gconsym\ \makebox{\tt )}$\\ 
$\it qop$\>\makebox[3.5em]{$\rightarrow$}$\it qconop$\\ 
$\it qconop$\>\makebox[3.5em]{$\rightarrow$}$\it gconsym$\\ 
$\it gconsym$\>\makebox[3.5em]{$\rightarrow$}$\it \makebox{\tt :}$
\end{tabbing}\end{flushleft}
\indexsyn{aexp}%

{\em Lists}\index{list} are written \mbox{$\it \makebox{\tt [}e_1\makebox{\tt ,}\ \ldots \makebox{\tt ,}\ e_k\makebox{\tt ]}$}, where
\mbox{$\it k\geq 1$}.  The list constructor is \mbox{\tt :}, and the empty list is denoted \mbox{\tt []}.
Standard operations on
lists are given in the Prelude (see Section~\ref{basic-lists}, and
Chapter~\ref{stdprelude} notably Section~\ref{preludelist}).

\outline{
\paragraph*{Translation:}  
The following identity holds:
\begin{center}
\bt{lcl}%
\struthack{17pt}%
\mbox{$\it \makebox{\tt [}e_1\makebox{\tt ,}\ \ldots \makebox{\tt ,}\ e_k\makebox{\tt ]}$}  & \mbox{$\it =$} & \mbox{$\it e_1\ \makebox{\tt :\ (}e_2\ \makebox{\tt :\ (}\ \ldots \ \makebox{\tt (}e_k\ \makebox{\tt :\ [])))}$}
\et
\end{center}
where \mbox{\tt :} and \mbox{\tt []} are constructors for lists, as defined in
the Prelude (see Section~\ref{basic-lists}).  The types
of \mbox{$\it e_1$} through \mbox{$\it e_k$} must all be the same (call it \mbox{$\it t\/$}), and the
type of the overall expression is \mbox{\tt [}\mbox{$\it t$}\mbox{\tt ]} (see Section~\ref{type-syntax}).
}
The constructor ``\mbox{\tt :}'' is reserved solely for list construction; like
\mbox{\tt []}, it is considered part of the language syntax, and cannot be hidden or redefined.
It is a right-associative operator, with precedence level 5 (Section~\ref{fixity}).

\section{Tuples}
\label{tuples}
%
\begin{flushleft}\it\begin{tabbing}
\hspace{0.6in}\=\hspace{3.1in}\=\kill
$\it aexp$\>\makebox[3.5em]{$\rightarrow$}$\it \makebox{\tt (}\ exp_1\ \makebox{\tt ,}\ \ldots \ \makebox{\tt ,}\ exp_k\ \makebox{\tt )}$\>\makebox[3em]{}$\it (k\geq 2)$\\ 
$\it $\>\makebox[3.5em]{$|$}$\it qcon$\\ 
$\it qcon$\>\makebox[3.5em]{$\rightarrow$}$\it \makebox{\tt (,}\{\makebox{\tt ,}\}\makebox{\tt )}$\\ 
$\it $
\end{tabbing}\end{flushleft}
\indexsyn{aexp}%

{\em Tuples}\index{tuple} are written \mbox{$\it \makebox{\tt (}e_1\makebox{\tt ,}\ \ldots \makebox{\tt ,}\ e_k\makebox{\tt )}$}, and may be
of arbitrary length \mbox{$\it k\geq 2$}.  The constructor for an \mbox{$\it n$}-tuple is denoted by
\mbox{\tt (,}\ldots\mbox{\tt ,)}, where there are \mbox{$\it n-1$} commas.  Thus \mbox{\tt (a,b,c)} and
\mbox{\tt (,,)\ a\ b\ c} denote the same value.
Standard operations on tuples are given
in the Prelude (see Section~\ref{basic-tuples} and Chapter~\ref{stdprelude}).

\outline{
\paragraph*{Translation:}  
\mbox{$\it \makebox{\tt (}e_1\makebox{\tt ,}\ \ldots \makebox{\tt ,}\ e_k\makebox{\tt )}$} for \mbox{$\it k\geq2$} is an instance of a \mbox{$\it k$}-tuple as
defined in the Prelude, and requires no translation.  If
\mbox{$\it t_1$} through \mbox{$\it t_k$} are the types of \mbox{$\it e_1$} through \mbox{$\it e_k$},
respectively, then the type of the resulting tuple is 
\mbox{$\it \makebox{\tt (}t_1\makebox{\tt ,}\ \ldots \makebox{\tt ,}\ t_k\makebox{\tt )}$} (see Section~\ref{type-syntax}).
}

\section{Unit Expressions and Parenthesized Expressions}
\label{unit-expression}
\index{unit expression}
%
\begin{flushleft}\it\begin{tabbing}
\hspace{0.6in}\=\hspace{3.1in}\=\kill
$\it aexp$\>\makebox[3.5em]{$\rightarrow$}$\it gcon$\\ 
$\it $\>\makebox[3.5em]{$|$}$\it \makebox{\tt (}\ exp\ \makebox{\tt )}$\\ 
$\it gcon$\>\makebox[3.5em]{$\rightarrow$}$\it \makebox{\tt ()}$
\end{tabbing}\end{flushleft}
\indexsyn{aexp}%

\noindent
The form \mbox{$\it \makebox{\tt (}e\makebox{\tt )}$} is simply a {\em parenthesized expression}, and is
equivalent to \mbox{$\it e$}.  The {\em unit expression} \mbox{\tt ()} has type
\mbox{\tt ()}\index{trivial type} (see
Section~\ref{type-syntax}).  It is the only member of that type apart
from $\bot$, and can
be thought of as the ``nullary tuple'' (see Section~\ref{basic-trivial}).
\nopagebreak[4]

\outline{
\paragraph{Translation:}  
\mbox{$\it \makebox{\tt (}e\makebox{\tt )}$} is equivalent to \mbox{$\it e$}.
}

\section{Arithmetic Sequences}
\label{arithmetic-sequences}
%
\begin{flushleft}\it\begin{tabbing}
\hspace{0.6in}\=\hspace{3.1in}\=\kill
$\it aexp$\>\makebox[3.5em]{$\rightarrow$}$\it \makebox{\tt [}\ exp_1\ [\makebox{\tt ,}\ exp_2]\ \makebox{\tt ..}\ [exp_3]\ \makebox{\tt ]}$
\end{tabbing}\end{flushleft}
\indexsyn{aexp}%

\noindent


The {\em arithmetic sequence}\index{arithmetic sequence}
\mbox{$\it \makebox{\tt [}e_1\makebox{\tt ,}\ e_2\ \makebox{\tt ..}\ e_3\makebox{\tt ]}$} denotes a list of values of
type \mbox{$\it t$}, where each of the \mbox{$\it e_i$} has type \mbox{$\it t$}, and \mbox{$\it t$} is an
instance of class \mbox{\tt Enum}.

\outline{
\paragraph{Translation:}
Arithmetic sequences satisfy these identities:
\begin{center}
\begin{tabular}{lcl}%
\struthack{17pt}%
\mbox{\tt [\ }\mbox{$\it e_1$}\mbox{\tt ..\ ]}         & \mbox{$\it =$} 
                        & \mbox{\tt enumFrom} \mbox{$\it e_1$} \\
\mbox{\tt [\ }\mbox{$\it e_1$}\mbox{\tt ,}\mbox{$\it e_2$}\mbox{\tt ..\ ]} & \mbox{$\it =$} 
                        & \mbox{\tt enumFromThen} \mbox{$\it e_1$} \mbox{$\it e_2$} \\
\mbox{\tt [\ }\mbox{$\it e_1$}\mbox{\tt ..}\mbox{$\it e_3$}\mbox{\tt \ ]}  & \mbox{$\it =$} 
                        & \mbox{\tt enumFromTo} \mbox{$\it e_1$} \mbox{$\it e_3$} \\
\mbox{\tt [\ }\mbox{$\it e_1$}\mbox{\tt ,}\mbox{$\it e_2$}\mbox{\tt ..}\mbox{$\it e_3$}\mbox{\tt \ ]} 
                        & \mbox{$\it =$} 
                        & \mbox{\tt enumFromThenTo} \mbox{$\it e_1$} \mbox{$\it e_2$} \mbox{$\it e_3$}
\end{tabular}
\end{center}
where \mbox{\tt enumFrom}, \mbox{\tt enumFromThen}, \mbox{\tt enumFromTo}, and \mbox{\tt enumFromThenTo}
are class methods in the class \mbox{\tt Enum} as defined in the Prelude
(see Figure~\ref{standard-classes}).
}

The semantics of arithmetic sequences therefore depends entirely
on the instance declaration for the type \mbox{$\it t$}.  
See Section~\ref{enum-class} for more details of which \mbox{\tt Prelude}
types are in \mbox{\tt Enum} and their semantics.

\section{List Comprehensions}
\index{list comprehension}
\index{let expression!in list comprehensions}
\label{list-comprehensions}
%
\begin{flushleft}\it\begin{tabbing}
\hspace{0.6in}\=\hspace{3.1in}\=\kill
$\it aexp$\>\makebox[3.5em]{$\rightarrow$}$\it \makebox{\tt [}\ exp\ \makebox{\tt |}\ qual_1\ \makebox{\tt ,}\ \ldots \ \makebox{\tt ,}\ qual_n\ \makebox{\tt ]}$\>\makebox[3em]{}$\it (\tr{list\ comprehension},\ n\geq 1)$\\ 
$\it qual$\>\makebox[3.5em]{$\rightarrow$}$\it pat\ \makebox{\tt <-}\ exp$\>\makebox[3em]{}$\it (\tr{generator})$\\ 
$\it $\>\makebox[3.5em]{$|$}$\it \makebox{\tt let}\ decls$\>\makebox[3em]{}$\it (\tr{local\ declaration})$\\ 
$\it $\>\makebox[3.5em]{$|$}$\it exp$\>\makebox[3em]{}$\it (\tr{boolean\ guard})$
\end{tabbing}\end{flushleft}
\indexsyn{aexp}
\indexsyn{qual}

\noindent
A {\em list comprehension} has the form \mbox{$\it \makebox{\tt [}\ e\ \makebox{\tt |}\ q_1\makebox{\tt ,}\ \ldots \makebox{\tt ,}\ q_n\ \makebox{\tt ]},\\
\it n\geq 1,$} where the \mbox{$\it q_i$} qualifiers\index{qualifier} are either
\begin{itemize}
\item {\em generators}\index{generator} of the form \mbox{$\it p\ \makebox{\tt <-}\ e$}, where
\mbox{$\it p$} is a
pattern (see Section~\ref{pattern-matching}) of type \mbox{$\it t$} and \mbox{$\it e$} is an
expression of type \mbox{$\it \makebox{\tt [}t\makebox{\tt ]}$}
\item {\em local bindings} that provide new definitions for use in
the generated expression \mbox{$\it e$} or subsequent boolean guards and generators
\item {\em boolean guards},\index{boolean guard} which are arbitrary expressions of
type \mbox{\tt Bool}.
\end{itemize}

Such a list comprehension returns the list of elements
produced by evaluating \mbox{$\it e$} in the successive environments
created by the nested, depth-first evaluation of the generators in the
qualifier list.  Binding of variables occurs according to the normal
pattern matching rules (see Section~\ref{pattern-matching}), and if a
match fails then that element of the list is simply skipped over.  Thus:\nopagebreak[4]
\bprog
{\small\begin{verbatim}
[ x |  xs   <- [ [(1,2),(3,4)], [(5,4),(3,2)] ], 
      (3,x) <- xs ]
\end{verbatim}}
\eprog
yields the list \mbox{\tt [4,2]}.  If a qualifier is a boolean guard, it must evaluate
to \mbox{\tt True} for the previous pattern match to succeed.  
As usual, bindings in list comprehensions can shadow those in outer
scopes; for example:
\[\ba{lll}
\mbox{\tt [\ x\ |\ x\ <-\ x,\ x\ <-\ x\ ]} & = & \mbox{\tt [\ z\ |\ y\ <-\ x,\ z\ <-\ y]} \\
\ea\]
\outline{
\paragraph{Translation:}
List comprehensions satisfy these identities, which may be
used as a translation into the kernel:
\begin{center}
\bt{lcl}%
\struthack{17pt}%
\mbox{$\it \makebox{\tt [\ }\ e\makebox{\tt \ |\ True\ ]}$} & = & \mbox{$\it \makebox{\tt [}e\makebox{\tt ]}$} \\
\mbox{$\it \makebox{\tt [\ }\ e\makebox{\tt \ |\ }q\makebox{\tt \ ]}$} & = & \mbox{$\it \makebox{\tt [}~\ e~\makebox{\tt |}~q\makebox{\tt ,\ True\ ]}$} \\
\mbox{$\it \makebox{\tt [\ }\ e\makebox{\tt \ |\ }b\makebox{\tt ,}~\ Q\ ~\makebox{\tt ]}$} & = &
        \mbox{$\it \makebox{\tt if}~b~\makebox{\tt then}~\makebox{\tt [\ }\ e\makebox{\tt \ |\ }Q\makebox{\tt \ ]}~\makebox{\tt else\ []}$} \\
\mbox{$\it \makebox{\tt [\ }\ e\makebox{\tt \ |\ }p\ \makebox{\tt <-}\ l\makebox{\tt ,}~\ Q\makebox{\tt \ ]}$} & = &
        \mbox{$\it \makebox{\tt let\ ok}~p~\makebox{\tt =}~\makebox{\tt [\ }\ e\makebox{\tt \ |\ }Q\makebox{\tt \ ]}$} \\
&&       \mbox{\tt \ \ \ \ ok\ {\char'137}\ =\ []} \\
&&      \mbox{$\it \makebox{\tt in\ concatMap\ ok}~\ l$} \\
\mbox{$\it \makebox{\tt [\ }\ e\makebox{\tt \ |\ let}~decls\makebox{\tt ,}~\ Q\makebox{\tt \ ]}$} & = &
        \mbox{$\it \makebox{\tt let}~decls~\makebox{\tt in}~\makebox{\tt [\ }\ e\makebox{\tt \ |\ }Q\makebox{\tt \ ]}$}
\et
\end{center}
where \mbox{$\it e$} ranges over expressions, \mbox{$\it p$} over
patterns, \mbox{$\it l$} over list-valued expressions, \mbox{$\it b$} over
boolean expressions, \mbox{$\it decls$} over declaration lists, 
\mbox{$\it q$} over qualifiers, and \mbox{$\it Q$} over sequences of qualifiers.  \mbox{$\it \makebox{\tt ok}$} is a fresh variable.
The function \mbox{\tt concatMap}, and boolean value \mbox{\tt True}, are defined in the Prelude.
}

As indicated by the translation of list comprehensions, variables
bound by \mbox{\tt let} have fully polymorphic types while those defined by
\mbox{\tt <-} are lambda bound and are thus monomorphic (see Section
\ref{monomorphism}).

\section{Let Expressions}
\index{let expression}
\label{let-expressions}
%
% Including this syntax blurb does REALLY bad things to page breaking
% in the 1.[12] report (sigh); ToDo: hope it goes away.
\begin{flushleft}\it\begin{tabbing}
\hspace{0.6in}\=\hspace{3.1in}\=\kill
$\it lexp$\>\makebox[3.5em]{$\rightarrow$}$\it \makebox{\tt let}\ decls\ \makebox{\tt in}\ exp$
\end{tabbing}\end{flushleft}
\indexsyn{exp}
\index{declaration!within a {\tt let} expression}

\noindent
{\em Let expressions} have the general form
\mbox{$\it \makebox{\tt let\ {\char'173}}\ d_1\ \makebox{\tt ;}\ \ldots \ \makebox{\tt ;}\ d_n\ \makebox{\tt {\char'175}\ in}\ e$},
and introduce a
nested, lexically-scoped, 
mutually-recursive list of declarations (\mbox{\tt let} is often called \mbox{\tt letrec} in
other languages).  The scope of the declarations is the expression \mbox{$\it e$}
and the right hand side of the declarations.  Declarations are
described in Chapter~\ref{declarations}.  Pattern bindings are matched
lazily; an implicit \mbox{\tt {\char'176}} makes these patterns
irrefutable.\index{irrefutable pattern}
For example, 

\bt{l}
\mbox{\tt let\ (x,y)\ =\ undefined\ in\ }\mbox{$\it e$}
\et

does not cause an execution-time error until \mbox{\tt x} or \mbox{\tt y} is evaluated.

\outline{\ifhtml{}{\small}
\paragraph*{Translation:} The dynamic semantics of the expression 
\mbox{$\it \makebox{\tt let\ {\char'173}}\ d_1\ \makebox{\tt ;}\ \ldots \ \makebox{\tt ;}\ d_n\ \makebox{\tt {\char'175}\ in}\ e_0$} are captured by this
translation: After removing all type signatures, each
declaration \mbox{$\it d_i$} is translated into an equation of the form 
\mbox{$\it p_i\ \makebox{\tt =}\ e_i$}, where \mbox{$\it p_i$} and \mbox{$\it e_i$} are patterns and expressions
respectively, using the translation in
Section~\ref{function-bindings}.  Once done, these identities
hold, which may be used as a translation into the kernel:
\begin{center}
\bt{lcl}%
\struthack{17pt}%
\mbox{\tt let\ {\char'173}}\mbox{$\it p_1$}\mbox{\tt =}\mbox{$\it e_1$}\mbox{\tt ;\ } ... \mbox{\tt ;\ }\mbox{$\it p_n$}\mbox{\tt =}\mbox{$\it e_n$}\mbox{\tt {\char'175}\ in} \mbox{$\it e_0$}
      &=& \mbox{\tt let\ ({\char'176}}\mbox{$\it p_1$}\mbox{\tt ,} ... \mbox{\tt ,{\char'176}}\mbox{$\it p_n$}\mbox{\tt )\ =\ (}\mbox{$\it e_1$}\mbox{\tt ,} ... \mbox{\tt ,}\mbox{$\it e_n$}\mbox{\tt )\ in} \mbox{$\it e_0$} \\
\mbox{\tt let\ }\mbox{$\it p$}\mbox{\tt \ =\ }\mbox{$\it e_1$} \mbox{\tt \ in\ } \mbox{$\it e_0$}
        &=& \mbox{\tt case\ }\mbox{$\it e_1$}\mbox{\tt \ of\ {\char'176}}\mbox{$\it p$}\mbox{\tt \ ->\ }\mbox{$\it e_0$}   \\
        & & {\rm where no variable in \mbox{$\it p$} appears free in \mbox{$\it e_1$}} \\
\mbox{\tt let\ }\mbox{$\it p$}\mbox{\tt \ =\ }\mbox{$\it e_1$} \mbox{\tt \ in\ } \mbox{$\it e_0$}
      &=& \mbox{\tt let\ }\mbox{$\it p$}\mbox{\tt \ =\ fix\ (\ {\char'134}\ {\char'176}}\mbox{$\it p$}\mbox{\tt \ ->\ }\mbox{$\it e_1$}\mbox{\tt )\ in} \mbox{$\it e_0$}
\et
\end{center}
where \mbox{\tt fix} is the least fixpoint operator.  Note the use of the
irrefutable patterns \mbox{$\it \makebox{\tt {\char'176}}p$}.  This translation
does not preserve the static semantics because the use of \mbox{\tt case}
precludes a fully polymorphic typing of the bound variables.
%\folks{Needs work -- KH}
% This same semantics applies to the top-level of
%a program that has been translated into a \mbox{\tt let} expression,
% as described at the beginning of Section~\ref{modules}.
The static semantics of the bindings in a \mbox{\tt let} expression
are described in 
Section~\ref{pattern-bindings}.
}

\section{Case Expressions}
\label{case}
%
\begin{flushleft}\it\begin{tabbing}
\hspace{0.6in}\=\hspace{3.1in}\=\kill
$\it lexp$\>\makebox[3.5em]{$\rightarrow$}$\it \makebox{\tt case}\ exp\ \makebox{\tt of}\ \makebox{\tt {\char'173}}\ alts\ \makebox{\tt {\char'175}}$\\ 
$\it alts$\>\makebox[3.5em]{$\rightarrow$}$\it alt_1\ \makebox{\tt ;}\ \ldots \ \makebox{\tt ;}\ alt_n$\>\makebox[3em]{}$\it (n\geq 1)$\\ 
$\it alt$\>\makebox[3.5em]{$\rightarrow$}$\it pat\ \makebox{\tt ->}\ exp\ [\makebox{\tt where}\ decls]$\\ 
$\it $\>\makebox[3.5em]{$|$}$\it pat\ gdpat\ [\makebox{\tt where}\ decls]$\\ 
$\it $\>\makebox[3.5em]{$|$}$\it $\>\makebox[3em]{}$\it (empty\ alternative)$\\ 
$\it $\\ 
$\it gdpat$\>\makebox[3.5em]{$\rightarrow$}$\it \hprime{guards}\ \makebox{\tt ->}\ exp\ [\ gdpat\ ]$\\ 
$\it \hprime{guards}$\>\makebox[3.5em]{$\rightarrow$}$\it \hprime{\makebox{\tt |}\ guard_1,\ \ldots ,\ guard_n}$\>\makebox[3em]{}$\it \hprime{(n\geq 1)}$\\ 
$\it \hprime{guard}$\>\makebox[3.5em]{$\rightarrow$}$\it \hprime{pat\ \makebox{\tt <-}\ infixexp}$\>\makebox[3em]{}$\it (\hprime{\tr{pattern\ guard}})$\\ 
$\it $\>\makebox[3.5em]{$|$}$\it \hprime{\makebox{\tt let}\ decls}$\>\makebox[3em]{}$\it (\hprime{\tr{local\ declaration}})$\\ 
$\it $\>\makebox[3.5em]{$|$}$\it infixexp$\>\makebox[3em]{}$\it (\tr{boolean\ guard})$
\end{tabbing}\end{flushleft}
\indexsyn{exp}%
\indexsyn{alts}%
\indexsyn{alt}%
\indexsyn{gdpat}%
\indexsyn{guards}%
\indexsyn{guard}%


A {\em case expression}\index{case expression} has the general form
\[
\mbox{$\it \makebox{\tt case}\ e\ \makebox{\tt of\ {\char'173}\ }p_1\ match_1\ \makebox{\tt ;}\ \ldots \ \makebox{\tt ;}\ p_n\ match_n\ \makebox{\tt {\char'175}}$}
\]
where each \mbox{$\it match_i$} is of the general form
\[\ba{lll}
 & \mbox{$\it \makebox{\tt |}\ \hprime{gs_{i1}}$}   & \mbox{$\it \makebox{\tt ->}\ e_{i1}$} \\
 & \mbox{$\it \ldots $} \\
 & \mbox{$\it \makebox{\tt |}\ \hprime{gs_{im_i}}$} & \mbox{$\it \makebox{\tt ->}\ e_{im_i}$} \\
 & \multicolumn{2}{l}{\mbox{$\it \makebox{\tt where}\ decls_i$}}
\ea\]
(Notice that in the syntax rule for \mbox{$\it guards$}, the ``\mbox{\tt |}'' is a 
terminal symbol, not the syntactic metasymbol for alternation.)
Each alternative \mbox{$\it p_i\ match_i$} consists of a 
pattern\index{pattern} \mbox{$\it p_i$} and its matches, \mbox{$\it match_i$}.
Each match in turn
consists of a sequence of pairs of guards\index{guard}
\mbox{$\it gs_{ij}$} and bodies \mbox{$\it e_{ij}$} (expressions), followed by
optional bindings (\mbox{$\it decls_i$}) that scope over all of the guards and
expressions of the alternative.

\begin{haskellprime}
\index{Pattern Guards}
\index{guards}
A {\em guard}\index{guard} has one of the following forms:
\begin{itemize}
\item {\em pattern guards}\index{pattern guard} are of the form \mbox{$\it p\ \makebox{\tt <-}\ e$}, where
\mbox{$\it p$} is a 
pattern (see Section~\ref{pattern-matching}) of type \mbox{$\it t$} and \mbox{$\it e$} is an
expression type \mbox{$\it t$}\footnote{Note that the syntax of a pattern guard is the same as that of a generator in a list comprehension. 
The contextual difference is that, in a list comprehension, a pattern of type \mbox{$\it t$} goes with an expression of type \mbox{$\it \makebox{\tt [}t\makebox{\tt ]}$}.}.
They succeed if the expression \mbox{$\it e$} matches the pattern \mbox{$\it p$}, and introduce the bindings of the pattern to the environment.
\item {\em local bindings} are of the form \mbox{$\it \makebox{\tt let}\ decls$}.  They always succeed, and they introduce the names defined in \mbox{$\it decls$} to the environment.
\item {\em boolean guards}\index{boolean guard} are arbitrary expressions of
type \mbox{\tt Bool}.  They succeed if the expression evaluates to \mbox{\tt True}, and they do not introduce new names to the environment.  A boolean guard, \mbox{$\it g$}, is semantically equivalent to the pattern guard \mbox{$\it \makebox{\tt True\ <-\ }g$}.
\end{itemize}
\end{haskellprime}


An alternative of the form
\[
\mbox{$\it pat\ \makebox{\tt ->}\ exp\ \makebox{\tt where}\ decls$}
\]
is treated as shorthand for:
\[\ba{lll}
 & \mbox{$\it pat\ \makebox{\tt |\ True}$}   & \mbox{$\it \makebox{\tt ->}\ exp$} \\
 & \multicolumn{2}{l}{\mbox{$\it \makebox{\tt where}\ decls$}}
\ea\]

A case expression must have at least one alternative and each alternative must
have at least one body.  Each body must have the same type, and the
type of the whole expression is that type.

A case expression is evaluated by pattern matching the expression \mbox{$\it e$}
against the individual alternatives.  The alternatives are tried
sequentially, from top to bottom.  If \mbox{$\it e$} matches the pattern of an
alternative, then the guarded expressions for that alternative are
tried sequentially from top to bottom in the environment of the case
expression extended first by the bindings created during the matching
of the pattern, and then by the \mbox{$\it decls_i$} in the \mbox{\tt where} clause
associated with that alternative.

\begin{haskellprime}
For each guarded expression, the comma-separated guards are tried
sequentially from left to right.  If all of them succeed, then the
corresponding expression is evaluated in the environment extended with
the bindings introduced by the guards.  That is, the bindings that are
introduced by a guard (either by using a let clause or a pattern
guard) are in scope in the following guards and the corresponding
expression.  If any of the guards fail, then this guarded expression
fails and the next guarded expression is tried.
\end{haskellprime}

If none of the guarded expressions for a given alternative succeed,
then matching continues with the next alternative.  If no alternative
succeeds, then the result is \mbox{$\it \bot$}.  Pattern matching is described in
Section~\ref{pattern-matching}, with the formal semantics of case
expressions in Section~\ref{case-semantics}.

{\em A note about parsing.} The expression
\bprog
{\small\begin{verbatim}
  case x of { (a,_) | let b = not a in b :: Bool -> a }
\end{verbatim}}
\eprog
is tricky to parse correctly.  It has a single unambiguous parse, namely
\bprog
{\small\begin{verbatim}
  case x of { (a,_) | (let b = not a in b :: Bool) -> a }
\end{verbatim}}
\eprog
However, the phrase \mbox{$\it \makebox{\tt Bool\ ->\ a}$} is syntactically valid as a type, and
parsers with limited lookahead may incorrectly commit to this choice, and hence
reject the program.  Programmers are advised, therefore, to avoid guards that
end with a type signature --- indeed that is why a \mbox{$\it guard$} contains 
an \mbox{$\it infixexp$} not an \mbox{$\it exp$}.

\section{Do Expressions}
\index{do expression}
\label{do-expressions}
\index{let expression!in do expressions}
\index{monad}
%
\begin{flushleft}\it\begin{tabbing}
\hspace{0.6in}\=\hspace{3.1in}\=\kill
$\it lexp$\>\makebox[3.5em]{$\rightarrow$}$\it \makebox{\tt do}\ \makebox{\tt {\char'173}}\ stmts\ \makebox{\tt {\char'175}}$\>\makebox[3em]{}$\it (\tr{do\ expression})$\\ 
$\it stmts$\>\makebox[3.5em]{$\rightarrow$}$\it stmt_1\ \ldots \ stmt_n\ exp\ [\makebox{\tt ;}]$\>\makebox[3em]{}$\it (n\geq 0)$\\ 
$\it stmt$\>\makebox[3.5em]{$\rightarrow$}$\it exp\ \makebox{\tt ;}$\\ 
$\it $\>\makebox[3.5em]{$|$}$\it pat\ \makebox{\tt <-}\ exp\ \makebox{\tt ;}$\\ 
$\it $\>\makebox[3.5em]{$|$}$\it \makebox{\tt let}\ decls\ \makebox{\tt ;}$\\ 
$\it $\>\makebox[3.5em]{$|$}$\it \makebox{\tt ;}$\>\makebox[3em]{}$\it (empty\ statement)$
\end{tabbing}\end{flushleft}
% The semicolons are done differently than for decls
%  Reason: to do it the same way would mean:
%       stmts -> stmt1 ; ... ; stmtn ; exp [;]
% Now, what happens if n=0?  Is there a ';' before the exp or not?
% Putting the ';' on the end of the stmt makes that clear.
\indexsyn{exp}%
\indexsyn{stmt}%
\indexsyn{stmts}%

A {\em do expression} provides a more conventional syntax for monadic programming.
It allows an expression such as 
\bprog
{\small\begin{verbatim}
  putStr "x: "    >> 
  getLine         >>= \l ->
  return (words l)
\end{verbatim}}
\eprog
to be written in a more traditional way as:
\bprog
{\small\begin{verbatim}
  do putStr "x: "
     l <- getLine
     return (words l)
\end{verbatim}}
\eprog
\outline{
\paragraph*{Translation:} 
Do expressions satisfy these identities, which may be
used as a translation into the kernel, after eliminating empty \mbox{$\it stmts$}:
\begin{center}
\bt{lcl}%
\struthack{17pt}%
\mbox{\tt do\ {\char'173}}\mbox{$\it e$}\mbox{\tt {\char'175}}                       &=& \mbox{$\it e$}\\
\mbox{\tt do\ {\char'173}}\mbox{$\it e$}\mbox{\tt ;}\mbox{$\it stmts$}\mbox{\tt {\char'175}}             &=& \mbox{$\it e$} \mbox{\tt >>\ do\ {\char'173}}\mbox{$\it stmts$}\mbox{\tt {\char'175}} \\
\mbox{\tt do\ {\char'173}}\mbox{$\it p$}\mbox{\tt \ <-\ }\mbox{$\it e$}\mbox{\tt ;\ }\mbox{$\it stmts$}\mbox{\tt {\char'175}}   &=& \mbox{\tt let\ ok\ }\mbox{$\it p$}\mbox{\tt \ =\ do\ {\char'173}}\mbox{$\it stmts$}\mbox{\tt {\char'175}}\\
                                   & & \mbox{\tt \ \ \ \ ok\ {\char'137}\ =\ fail\ "..."}\\
                                   & & \mbox{\tt \ \ in\ }\mbox{$\it e$}\mbox{\tt \ >>=\ ok} \\
\mbox{\tt do\ {\char'173}let} \mbox{$\it decls$}\mbox{\tt ;\ }\mbox{$\it stmts$}\mbox{\tt {\char'175}}  &=& \mbox{\tt let} \mbox{$\it decls$} \mbox{\tt in\ do\ {\char'173}}\mbox{$\it stmts$}\mbox{\tt {\char'175}}\\
\et
\end{center}
The ellipsis \mbox{\tt "..."} stands for a compiler-generated error message,
passed to \mbox{\tt fail}, preferably giving some indication of the location
of the pattern-match failure;
the functions \mbox{\tt >>}, \mbox{\tt >>=}, and \mbox{\tt fail} are operations in the class \mbox{\tt Monad},
as defined in the Prelude\indexclass{Monad}; and \mbox{\tt ok} is a fresh
identifier. 
}
As indicated by the translation of \mbox{\tt do}, variables bound by \mbox{\tt let} have
fully polymorphic types while those defined by \mbox{\tt <-} are lambda bound
and are thus monomorphic.

\section{Datatypes with Field Labels}
\label{field-ops}
\index{data declaration@{\tt data} declaration}
\index{label}
\index{field label|see{label}}
A datatype declaration may optionally define field labels
(see Section~\ref{datatype-decls}).
These field labels can be used to 
construct, select from, and update fields in a manner
that is independent of the overall structure of the datatype.

Different datatypes cannot share common field labels in the same scope.
A field label can be used at most once in a constructor.
Within a datatype, however, a field label can be used in more
than one constructor provided the field has the same typing in all
constructors. To illustrate the last point, consider:
\bprog
{\small\begin{verbatim}
  data S = S1 { x :: Int } | S2 { x :: Int }   -- OK
  data T = T1 { y :: Int } | T2 { y :: Bool }  -- BAD
\end{verbatim}}
\eprog
Here \mbox{\tt S} is legal but \mbox{\tt T} is not, because \mbox{\tt y} is given 
inconsistent typings in the latter.

\subsection{Field Selection}
\begin{flushleft}\it\begin{tabbing}
\hspace{0.6in}\=\hspace{3.1in}\=\kill
$\it aexp$\>\makebox[3.5em]{$\rightarrow$}$\it qvar$
\end{tabbing}\end{flushleft}
\index{field label!selection}

Field labels are used as selector functions.  When used as a variable,
a field label serves as a function that extracts the field from an
object.  Selectors are top level bindings and so they
may be shadowed by local variables but cannot conflict with 
other top level bindings of the same name.  This shadowing only
affects selector functions; in record construction (Section~\ref{record-construction}) 
and update (Section~\ref{record-update}), field labels
cannot be confused with ordinary variables. 

\outline{
\paragraph*{Translation:} 
A field label \mbox{$\it f$} introduces a selector function defined as:
\begin{center}
\bt{lcl}
\struthack{17pt}%
\mbox{$\it f$}\mbox{\tt \ x} &=&\mbox{\tt case\ x\ of\ {\char'173}} \mbox{$\it C_1\ p_{11}\ \ldots\ p_{1k}$} \mbox{\tt \ ->\ } \mbox{$\it e_1$} \mbox{\tt ;} 
 \mbox{$\it \ldots$} \mbox{\tt ;} \mbox{$\it C_n\ p_{n1}\ \ldots\ p_{nk}$} \mbox{\tt \ ->\ } \mbox{$\it e_n$} \mbox{\tt {\char'175}}\\
\et
\end{center}
where \mbox{$\it C_1\ \ldots\ C_n$} are all the constructors of the datatype containing a
field labeled with \mbox{$\it f$}, \mbox{$\it p_{ij}$} is \mbox{\tt y} when \mbox{$\it f$} labels the \mbox{$\it j$}th
component of \mbox{$\it C_i$} or \mbox{\tt {\char'137}} otherwise, and \mbox{$\it e_i$} is \mbox{\tt y} when some field in
\mbox{$\it C_i$} has a label of \mbox{$\it f$} or \mbox{\tt undefined} otherwise.
}
\subsection{Construction Using Field Labels}
\label{record-construction}
\index{field label!construction}
\begin{flushleft}\it\begin{tabbing}
\hspace{0.6in}\=\hspace{3.1in}\=\kill
$\it aexp$\>\makebox[3.5em]{$\rightarrow$}$\it qcon\ \makebox{\tt {\char'173}}\ fbind_1\ \makebox{\tt ,}\ \ldots \ \makebox{\tt ,}\ fbind_n\ \makebox{\tt {\char'175}}$\>\makebox[3em]{}$\it (\tr{labeled\ construction},\ n\geq 0)$\\ 
$\it fbind$\>\makebox[3.5em]{$\rightarrow$}$\it qvar\ \makebox{\tt =}\ exp$
\end{tabbing}\end{flushleft}
\indexsyn{fbind}%

A constructor with labeled fields may be used to construct a value 
in which the components are specified by name rather than by position.
Unlike the braces used in declaration lists, these are not subject to
layout; the \mbox{\tt {\char'173}} and \mbox{\tt {\char'175}} characters must be explicit.  (This is also
true of field updates and field patterns.)
Construction using field labels is subject to the following constraints:
\begin{itemize}
\item Only field labels declared with the specified constructor may be
mentioned. 
\item A field label may not be mentioned more than once.
\item Fields not mentioned are initialized to $\bot$.
\item A compile-time error occurs when any strict fields (fields
whose declared types are prefixed by \mbox{\tt !}) are omitted during
construction.  Strict fields are discussed in Section~\ref{strictness-flags}.
\end{itemize}
The expression \mbox{\tt F\ {\char'173}{\char'175}}, where \mbox{\tt F} is a data constructor, is legal 
{\em whether or not \mbox{\tt F} was declared with record syntax} (provided \mbox{\tt F} has no strict
fields --- see the fourth bullet above); 
it denotes \mbox{$\it \makebox{\tt F}\ \bot_1\ \ldots \ \bot_n$}, where \mbox{$\it n$} is the arity of \mbox{\tt F}.

\outline{
\paragraph*{Translation:} 
In the binding \mbox{$\it f$} \mbox{\tt =} \mbox{$\it v$}, the field \mbox{$\it f$} labels \mbox{$\it v$}.
\begin{center}
\bt{lcl}
\struthack{17pt}%
\mbox{$\it C$} \mbox{\tt {\char'173}} \mbox{$\it bs$} \mbox{\tt {\char'175}} &=& \mbox{$\it C\ (pick^C_1\ bs\ \makebox{\tt undefined})\ \ldots\ (pick^C_k\ bs\ \makebox{\tt undefined})$}\\
\et
\end{center}
where \mbox{$\it k$} is the arity of \mbox{$\it C$}.

The auxiliary function \mbox{$\it pick^C_i\ bs\ d$} is defined as follows:
\begin{quote}
    If the \mbox{$\it i$}th component of a constructor \mbox{$\it C$} has the
    field label \mbox{$\it f$}, and if \mbox{$\it f=v$} appears in the binding list
    \mbox{$\it bs$}, then \mbox{$\it pick^C_i\ bs\ d$} is \mbox{$\it v$}.  Otherwise, \mbox{$\it pick^C_i\ bs\ d$} is
    the default value \mbox{$\it d$}.
\end{quote}
}

\subsection{Updates Using Field Labels}
\label{record-update}
\index{field label!update}
\begin{flushleft}\it\begin{tabbing}
\hspace{0.6in}\=\hspace{3.1in}\=\kill
$\it aexp$\>\makebox[3.5em]{$\rightarrow$}$\it aexp_{\langle{}qcon\rangle{}}\ \makebox{\tt {\char'173}}\ fbind_1\ \makebox{\tt ,}\ \ldots \ \makebox{\tt ,}\ fbind_n\ \makebox{\tt {\char'175}}$\>\makebox[3em]{}$\it (\tr{labeled\ update},\ n\geq 1)$
\end{tabbing}\end{flushleft}

Values belonging to a datatype with field labels may be
non-destructively updated.  This creates a new value in which the
specified field values replace those in the existing value.  
Updates are restricted in the following ways:
\begin{itemize}
\item All labels must be taken from the same datatype.
\item At least one constructor must define all of the labels
mentioned in the update.
\item No label may be mentioned more than once.
\item An execution error occurs when the value being updated does
not contain all of the specified labels.
\end{itemize}
\outline{
\paragraph*{Translation:} 
Using the prior definition of \mbox{$\it pick$},
\begin{center}
\bt{lcl}
\struthack{17pt}%
\mbox{$\it e$} \mbox{\tt {\char'173}} \mbox{$\it bs$} \mbox{\tt {\char'175}} &=& \mbox{\tt case} \mbox{$\it e$} \mbox{\tt of}\\
&&\mbox{\tt \ \ \ \ \ \ \ \ }\mbox{$\it C_1\ v_1\ \ldots \ v_{k_1}$} \mbox{\tt ->} \mbox{$\it C_1\ (pick^{C_1}_1\ bs\ v_1)\ \ldots \ (pick^{C_1}_{k_1}\ bs\ v_{k_1})$}\\
&&\mbox{\tt \ \ \ \ \ \ \ \ \ \ \ \ } ... \\
&&\mbox{\tt \ \ \ \ \ \ \ \ }\mbox{$\it C_j\ v_1\ \ldots \ v_{k_j}$} \mbox{\tt ->} \mbox{$\it C_j\ (pick^{C_j}_1\ bs\ v_1)\ \ldots \ (pick^{C_j}_{k_j}\ bs\ v_{k_j})$}\\
&&\mbox{\tt \ \ \ \ \ \ \ \ {\char'137}\ ->\ error\ "Update\ error"}\\
\et
\end{center}
where \mbox{$\it \{C_1,\ldots ,C_j\}$} is the set of constructors containing all labels
in \mbox{$\it bs$}, and \mbox{$\it k_i$} is the arity of \mbox{$\it C_i$}.
}
Here are some examples using labeled fields:
\bprog
{\small\begin{verbatim}
data T    = C1 {f1,f2 :: Int}
          | C2 {f1 :: Int,
                f3,f4 :: Char}
\end{verbatim}}
\eprog
\[\bt{|l|l|}%%b
\hline
Expression                                  & Translation                       \\
\hline
\mbox{\tt C1\ {\char'173}f1\ =\ 3{\char'175}}                       & \mbox{\tt C1\ 3\ undefined}          \\
\mbox{\tt C2\ {\char'173}f1\ =\ 1,\ f4\ =\ 'A',\ f3\ =\ 'B'{\char'175}}   & \mbox{\tt C2\ 1\ 'B'\ 'A'}            \\
\mbox{\tt x\ {\char'173}f1\ =\ 1{\char'175}}                 & \mbox{\tt case\ x\ of\ C1\ {\char'137}\ f2\ \ \ \ ->\ C1\ 1\ f2} \\
                             & \mbox{\tt \ \ \ \ \ \ \ \ \ \ C2\ {\char'137}\ f3\ f4\ ->\ C2\ 1\ f3\ f4}   \\
\hline\et\]
The field \mbox{\tt f1} is common to both constructors in T.  This
example translates expressions using constructors in field-label
notation into equivalent expressions using the same constructors
without field labels. 
A compile-time error will result if no single constructor
defines the set of field labels used in an update, such as
\mbox{\tt x\ {\char'173}f2\ =\ 1,\ f3\ =\ 'x'{\char'175}}. 

\section{Expression Type-Signatures}
\index{expression type-signature}
\label{expression-type-sigs}
%
\begin{flushleft}\it\begin{tabbing}
\hspace{0.6in}\=\hspace{3.1in}\=\kill
$\it exp$\>\makebox[3.5em]{$\rightarrow$}$\it exp\ \makebox{\tt ::}\ [context\ \makebox{\tt =>}]\ type$
\end{tabbing}\end{flushleft}
\indexsyn{exp}
\indextt{::}

\nopagebreak[4]
{\em Expression type-signatures} have the form \mbox{$\it e\ \makebox{\tt ::}\ t$}, where \mbox{$\it e$}
is an expression and \mbox{$\it t$} is a type (Section~\ref{type-syntax}); they
are used to type an expression explicitly
and may be used to resolve ambiguous typings due to overloading (see
Section~\ref{default-decls}).  The value of the expression is just that of
\mbox{$\it exp$}.  As with normal type signatures (see
Section~\ref{type-signatures}), the declared type may be more specific than 
the principal type derivable from \mbox{$\it exp$}, but it is an error to give
a type that is more general than, or not comparable to, the
principal type.
\outline{
\paragraph*{Translation:} 
\begin{center}
\bt{lcl}
\struthack{17pt}%
\mbox{$\it e\ \makebox{\tt ::}\ t$} & = & \mbox{$\it \makebox{\tt let\ {\char'173}}\ v\ \makebox{\tt ::}\ t\makebox{\tt ;\ }\ v\ \makebox{\tt =}\ e\ \makebox{\tt {\char'175}\ in\ }v$}
\et
\end{center}
}


\section{Pattern Matching}
\index{pattern-matching}
\label{pattern-matching}
\label{patterns}

{\em Patterns} appear in lambda abstractions, function definitions, pattern
bindings, list comprehensions, do expressions, and case expressions.
However, the 
first five of these ultimately translate into case expressions, so
defining the semantics of pattern matching for case expressions is sufficient.
%it suffices to restrict the definition of the semantics of
%pattern-matching to case expressions.

\subsection{Patterns}
\label{pattern-definitions}

Patterns\index{pattern} have this syntax:
\begin{flushleft}\it\begin{tabbing}
\hspace{0.6in}\=\hspace{3.1in}\=\kill
$\it pat$\>\makebox[3.5em]{$\rightarrow$}$\it \hprime{lpat\ qconop\ pat}$\>\makebox[3em]{}$\it (\tr{infix\ constructor})$\\ 
$\it $\>\makebox[3.5em]{$|$}$\it \hprime{lpat}$\\ 
$\it $\\ 
$\it \hprime{lpat}$\>\makebox[3.5em]{$\rightarrow$}$\it apat$\\ 
$\it $\>\makebox[3.5em]{$|$}$\it \hprime{\makebox{\tt -}\ (integer\ |\ float)}$\>\makebox[3em]{}$\it (\tr{negative\ literal})$\\ 
$\it $\>\makebox[3.5em]{$|$}$\it gcon\ apat_1\ \ldots \ apat_k$\>\makebox[3em]{}$\it (\tr{arity}\ gcon\ =\ k,\ k\geq 1)$\\ 
$\it $\\ 
$\it apat$\>\makebox[3.5em]{$\rightarrow$}$\it var\ [{\tt\ @}\ apat]$\>\makebox[3em]{}$\it (\tr{as\ pattern})$\\ 
$\it $\>\makebox[3.5em]{$|$}$\it gcon$\>\makebox[3em]{}$\it (\tr{arity}\ gcon\ =\ 0)$\\ 
$\it $\>\makebox[3.5em]{$|$}$\it qcon\ \makebox{\tt {\char'173}}\ fpat_1\ \makebox{\tt ,}\ \ldots \ \makebox{\tt ,}\ fpat_k\ \makebox{\tt {\char'175}}$\>\makebox[3em]{}$\it (\tr{labeled\ pattern},\ k\geq 0)$\\ 
$\it $\>\makebox[3.5em]{$|$}$\it literal$\\ 
$\it $\>\makebox[3.5em]{$|$}$\it \makebox{\tt {\char'137}}$\>\makebox[3em]{}$\it (\tr{wildcard})$\\ 
$\it $\>\makebox[3.5em]{$|$}$\it \makebox{\tt (}\ pat\ \makebox{\tt )}$\>\makebox[3em]{}$\it (\tr{parenthesized\ pattern})$\\ 
$\it $\>\makebox[3.5em]{$|$}$\it \makebox{\tt (}\ pat_1\ \makebox{\tt ,}\ \ldots \ \makebox{\tt ,}\ pat_k\ \makebox{\tt )}$\>\makebox[3em]{}$\it (\tr{tuple\ pattern},\ k\geq 2)$\\ 
$\it $\>\makebox[3.5em]{$|$}$\it \makebox{\tt [}\ pat_1\ \makebox{\tt ,}\ \ldots \ \makebox{\tt ,}\ pat_k\ \makebox{\tt ]}$\>\makebox[3em]{}$\it (\tr{list\ pattern},\ k\geq 1)$\\ 
$\it $\>\makebox[3.5em]{$|$}$\it \makebox{\tt {\char'176}}\ apat$\>\makebox[3em]{}$\it (\tr{irrefutable\ pattern})$\\ 
$\it $\\ 
$\it fpat$\>\makebox[3.5em]{$\rightarrow$}$\it qvar\ \makebox{\tt =}\ pat$
\end{tabbing}\end{flushleft}
\indexsyn{pat}%
\indexsyn{lpat}%
\indexsyn{apat}%
\indexsyn{fpats}%
\indexsyn{fpat}%

The arity of a constructor must match the number of
sub-patterns associated with it; one cannot match against a
partially-applied constructor.

All patterns must be {\em linear}\index{linearity}
\index{linear pattern}---no variable may appear more than once.  For
example, this definition is illegal:
\bprog
{\small\begin{verbatim}
f (x,x) = x     -- ILLEGAL; x used twice in pattern
\end{verbatim}}
\eprog
Patterns of the form \mbox{$\it var$}{\tt @}\mbox{$\it pat$} are called {\em as-patterns},
\index{as-pattern ({\tt {\char'100}})}
and allow one to use \mbox{$\it var$}
as a name for the value being matched by \mbox{$\it pat$}.  For example,\nopagebreak[4]
\bprog
{\small\begin{verbatim}
case e of { xs@(x:rest) -> if x==0 then rest else xs }
\end{verbatim}}
\eprog
is equivalent to:
\bprog
{\small\begin{verbatim}
let { xs = e } in
  case xs of { (x:rest) -> if x==0 then rest else xs }
\end{verbatim}}
\eprog

Patterns of the form \mbox{\tt {\char'137}} are {\em
wildcards}\index{wildcard pattern ({\tt {\char'137}})} and are useful when some part of a pattern
is not referenced on the right-hand-side.  It is as if an
identifier not used elsewhere were put in its place.  For example,
\bprog
{\small\begin{verbatim}
case e of { [x,_,_]  ->  if x==0 then True else False }
\end{verbatim}}
\eprog
is equivalent to:
\bprog
{\small\begin{verbatim}
case e of { [x,y,z]  ->  if x==0 then True else False }
\end{verbatim}}
\eprog
% where \mbox{\tt y} and \mbox{\tt z} are identifiers not used elsewhere.

%old:
%This translation is also 
%assumed prior to the semantics given below.

\subsection{Informal Semantics of Pattern Matching}

Patterns are matched against values.  Attempting to match a pattern
can have one of three results: it may {\em fail\/}; it may {\em
succeed}, returning a binding for each variable in the pattern; or it
may {\em diverge} (i.e.~return \mbox{$\it \bot$}).  Pattern matching proceeds
from left to right, and outside to inside, according to the following rules:
\begin{enumerate}
\item Matching the pattern \mbox{$\it var$} 
against a value \mbox{$\it v$} always succeeds and binds \mbox{$\it var$} to \mbox{$\it v$}.

\item
Matching the pattern \mbox{$\it \makebox{\tt {\char'176}}apat$} against a value \mbox{$\it v$} always succeeds.
The free variables in \mbox{$\it apat$} are bound to the appropriate values if matching
\mbox{$\it apat$} against \mbox{$\it v$} would otherwise succeed, and to \mbox{$\it \bot$} if matching
\mbox{$\it apat$} against \mbox{$\it v$} fails or diverges.  (Binding does {\em
not} imply evaluation.)

Operationally, this means that no matching is done on a
\mbox{$\it \makebox{\tt {\char'176}}apat$} pattern until one of the variables in \mbox{$\it apat$} is used.
At that point the entire pattern is matched against the value, and if
the match fails or diverges, so does the overall computation.

\item
Matching the wildcard pattern \mbox{\tt {\char'137}} against any value always succeeds,
and no binding is done.

\item
Matching the pattern \mbox{$\it con\ pat$} against a value, where \mbox{$\it con$} is a
constructor defined by \mbox{\tt newtype}, depends on the value:
\begin{itemize}
\item If the value is of the form \mbox{$\it con\ v$}, then \mbox{$\it pat$} is matched against \mbox{$\it v$}.
\item If the value is \mbox{$\it \bot$}, then \mbox{$\it pat$} is matched against \mbox{$\it \bot$}.
\end{itemize}
That is, constructors associated with
\mbox{\tt newtype} serve only to change the type of a value.\index{newtype declaration@{\tt newtype} declaration}

\item
Matching the pattern \mbox{$\it con\ pat_1\ \ldots\ pat_n$} against a value, where \mbox{$\it con$} is a
constructor defined by \mbox{\tt data}, depends on the value:
\begin{itemize}
\item If the value is of the form \mbox{$\it con\ v_1\ \ldots\ v_n$}, 
sub-patterns are matched left-to-right against the components of the data value;
if all matches succeed, the overall match
succeeds; the first to fail or diverge causes the overall match to
fail or diverge, respectively.  

\item If the value is of the form \mbox{$\it con'\ v_1\ \ldots\ v_m$}, where \mbox{$\it con$} is a different 
constructor to \mbox{$\it con'$}, the match fails.

\item If the value is \mbox{$\it \bot$}, the match diverges.
\end{itemize}

\item 
Matching against a constructor using labeled fields is the same as
matching ordinary constructor patterns except that the fields are
matched in the order they are named in the field list.  All fields
listed must be declared by the constructor; fields may not be named
more than once.  Fields not named by the pattern are ignored (matched
against \mbox{\tt {\char'137}}).

\item Matching a numeric, character, or string literal pattern \mbox{$\it k$} against a value \mbox{$\it v$}
\index{literal pattern}
succeeds if \mbox{$\it v\ ~\makebox{\tt ==}\ ~k$}, where \mbox{\tt ==}
is overloaded based on the type of the pattern.  The match diverges if
this test diverges.

The interpretation of numeric literals is exactly as described in Section~\ref{vars-and-lits};
that is, the overloaded function \mbox{\tt fromInteger} or \mbox{\tt fromRational} is 
applied to an \mbox{\tt Integer} or \mbox{\tt Rational} literal (resp)
to convert it to the appropriate type.

\item
Matching an as-pattern \mbox{$\it var$}{\tt @}\mbox{$\it apat$} against a value \mbox{$\it v$} is
\index{as-pattern ({\tt {\char'100}})}
the result of matching \mbox{$\it apat$} against \mbox{$\it v$}, augmented with the binding of
\mbox{$\it var$} to \mbox{$\it v$}.  If the match of \mbox{$\it apat$} against \mbox{$\it v$} fails or diverges,
then so does the overall match.
\end{enumerate}

Aside from the obvious static type constraints (for
example, it is a static error to match a character against a
boolean), the following static class constraints hold: 
\begin{itemize}
\item An integer
literal pattern
\index{integer literal pattern}
can only be matched against a value in the class
\mbox{\tt Num}.
\item A floating literal pattern
\index{floating literal pattern}
can only be matched against a value
in the class \mbox{\tt Fractional}.
\end{itemize}

It is sometimes helpful to distinguish two kinds of
patterns.  Matching an {\em irrefutable pattern}
\index{irrefutable pattern}
is non-strict: the pattern matches even if the value to be matched is \mbox{$\it \bot$}.
Matching a {\em refutable} pattern is strict: if the value to be matched
\index{refutable pattern}
is \mbox{$\it \bot$} the match diverges.
The irrefutable patterns are as follows:
a variable, a wildcard, \mbox{$\it N\ apat$} where \mbox{$\it N$} is a constructor
defined by \mbox{\tt newtype} and \mbox{$\it apat$} is irrefutable (see
Section~\ref{datatype-renaming}), 
\index{newtype declaration@{\tt newtype} declaration} 
\mbox{$\it var$}{\tt @}\mbox{$\it apat$} where \mbox{$\it apat$} is irrefutable,
or of the form \mbox{$\it \makebox{\tt {\char'176}}apat$} (whether or not \mbox{$\it apat$} is irrefutable).
All other patterns are {\em refutable}.

Here are some examples:
\begin{enumerate}
\item If the pattern \mbox{\tt ['a','b']} is matched against \mbox{$\it \makebox{\tt ['x',}\bot\makebox{\tt ]}$}, then \mbox{\tt 'a'}
\mbox{$\it fails$} to match against \mbox{\tt 'x'}, and the result is a failed match.  But
if \mbox{\tt ['a','b']} is matched against \mbox{$\it \makebox{\tt [}\bot\makebox{\tt ,'x']}$}, then attempting to match
\mbox{\tt 'a'} against \mbox{$\it \bot$} causes the match to \mbox{$\it diverge$}.

\item These examples demonstrate refutable vs.~irrefutable
matching:

\begin{tabular}{l}
\mbox{\tt ({\char'134}\ {\char'176}(x,y)\ ->\ 0)\ }$\bot$\mbox{\tt \ \ \ \ }$\Rightarrow$\mbox{\tt \ \ \ \ 0}\\
\mbox{\tt ({\char'134}\ \ (x,y)\ ->\ 0)\ }$\bot$\mbox{\tt \ \ \ \ }$\Rightarrow$\mbox{\tt \ \ \ \ }$\bot$
\end{tabular}

\begin{tabular}{l}
\mbox{\tt ({\char'134}\ {\char'176}[x]\ ->\ 0)\ []\ \ \ \ }$\Rightarrow$\mbox{\tt \ \ \ \ 0}\\
\mbox{\tt ({\char'134}\ {\char'176}[x]\ ->\ x)\ []\ \ \ \ }$\Rightarrow$\mbox{\tt \ \ \ \ }$\bot$
\end{tabular}

\begin{tabular}{l}
\mbox{\tt ({\char'134}\ {\char'176}[x,{\char'176}(a,b)]\ ->\ x)\ [(0,1),}$\bot$\mbox{\tt ]\ \ \ \ }$\Rightarrow$\mbox{\tt \ \ \ \ (0,1)}\\
\mbox{\tt ({\char'134}\ {\char'176}[x,\ (a,b)]\ ->\ x)\ [(0,1),}$\bot$\mbox{\tt ]\ \ \ \ }$\Rightarrow$\mbox{\tt \ \ \ \ }$\bot$
\end{tabular}

\begin{tabular}{l}
\mbox{\tt ({\char'134}\ \ (x:xs)\ ->\ x:x:xs)\ }$\bot$\mbox{\tt \ \ \ }$\Rightarrow$\mbox{\tt \ \ \ }$\bot$\\
\mbox{\tt ({\char'134}\ {\char'176}(x:xs)\ ->\ x:x:xs)\ }$\bot$\mbox{\tt \ \ \ }$\Rightarrow$\mbox{\tt \ \ \ }$\bot$\mbox{\tt :}$\bot$\mbox{\tt :}$\bot$
\end{tabular}

\item 
Consider the following declarations:
\bprog
{\small\begin{verbatim}
  newtype N = N Bool
  data    D = D !Bool
\end{verbatim}}
\eprog
These examples illustrate the difference in pattern matching
between types defined by \mbox{\tt data} and \mbox{\tt newtype}:

\begin{tabular}{l}
\mbox{\tt ({\char'134}\ \ (N\ True)\ ->\ True)\ }\mbox{$\it \bot$}\mbox{\tt \ \ \ \ \ }\mbox{$\it \Rightarrow$}\mbox{\tt \ \ \ \ }\mbox{$\it \bot$}\\
\mbox{\tt ({\char'134}\ \ (D\ True)\ ->\ True)\ }\mbox{$\it \bot$}\mbox{\tt \ \ \ \ \ }\mbox{$\it \Rightarrow$}\mbox{\tt \ \ \ \ }\mbox{$\it \bot$}\\
\mbox{\tt ({\char'134}\ {\char'176}(D\ True)\ ->\ True)\ }\mbox{$\it \bot$}\mbox{\tt \ \ \ \ \ }\mbox{$\it \Rightarrow$}\mbox{\tt \ \ \ \ True}
\end{tabular}

Additional examples may be found in Section~\ref{datatype-renaming}.

\end{enumerate}
\nopagebreak[4]
% \bprog
% @
% (\ t -> let (x,y) = t in 0) @\mbox{$\it \bot$}@    @\mbox{$\it \Rightarrow$}@    0
% (\ (x,y) -> 0) @\mbox{$\it \bot$}@    @\mbox{$\it \Rightarrow$}@    @\mbox{$\it \bot$}@
% @
% \eprog
% \bprog
% @
% (\ l -> let [x] = l in 0) []    @\mbox{$\it \Rightarrow$}@    0
% (\ l -> let [x] = l in x) []    @\mbox{$\it \Rightarrow$}@    @\mbox{$\it \bot$}@
% @
% \eprog
% \bprog
% @
% (\ l -> let [x,t] = l; (a,b) = t in x) [(0,1),@\mbox{$\it \bot$}@]    @\mbox{$\it \Rightarrow$}@    (0,1)
% (\ l -> let [x, (a,b)] = l in x) [(0,1),@\mbox{$\it \bot$}@]    @\mbox{$\it \Rightarrow$}@    @\mbox{$\it \bot$}@
% @
% \eprog
% \bprog
% @
% (\  (x:xs) -> x:x:xs) @\mbox{$\it \bot$}@   @\mbox{$\it \Rightarrow$}@   @\mbox{$\it \bot$}@
% (\ l -> let (x:xs) = l in x:x:xs) @\mbox{$\it \bot$}@   @\mbox{$\it \Rightarrow$}@   @\mbox{$\it \bot$}@:@\mbox{$\it \bot$}@:@\mbox{$\it \bot$}@
% @
% \eprogNoSkip
% \end{enumerate}

Top level patterns in case expressions and the set of top level
patterns in function or pattern bindings may have zero or more
associated {\em guards}\index{guards}.  See
Section~\ref{case} for the syntax and semantics of guards.

The guard semantics have an influence on the
strictness characteristics of a function or case expression.  In
particular, an otherwise irrefutable pattern
\index{irrefutable pattern}
may be evaluated because of a guard.  For example, in
\bprog
{\small\begin{verbatim}
f :: (Int,Int,Int) -> [Int] -> Int
f ~(x,y,z) [a] | (a == y) = 1
\end{verbatim}}
\eprog
% \bprog
% @
% f t [a] | a==y = 1 
%           where (x,y,z) = t
% @
% \eprog
both \mbox{\tt a} and \mbox{\tt y} will be evaluated by \mbox{\tt ==} in the guard.


\subsection{Formal Semantics of Pattern Matching}
\label{case-semantics}

The semantics of all pattern matching constructs other than \mbox{\tt case}
expressions are defined by giving identities that relate those
constructs to \mbox{\tt case} expressions.  The semantics of
\mbox{\tt case} expressions themselves are in turn given as a series of
identities, in Figures~\ref{simple-case-expr-1}--\ref{simple-case-expr-3}. 
Any implementation should behave so that these identities hold; it is 
not expected that it will use them directly, since that 
would generate rather inefficient code.
  
\begin{figure}[tb]
\outlinec{\small
\begin{tabular}{@{}cl}
(a)&\mbox{\tt case\ }$e$\mbox{\tt \ of\ {\char'173}\ }$alts$\mbox{\tt \ {\char'175}\ }$=$\mbox{\tt \ ({\char'134}}$v$\mbox{\tt \ ->\ case\ }$v$\mbox{\tt \ of\ {\char'173}\ }$alts$\mbox{\tt \ {\char'175})\ }\mbox{$\it e$}\\
&{\rm where $v$ is a new variable}\\
(b)&\mbox{\tt case\ }$v$\mbox{\tt \ of\ {\char'173}\ }$p_1\ \ match_1$\mbox{\tt ;\ \ }$\ldots{}$\mbox{\tt \ ;\ }$p_n\ \ match_n$\mbox{\tt \ {\char'175}}\\
   &$=$\mbox{\tt \ \ case\ }$v$\mbox{\tt \ of\ {\char'173}\ }$p_1\ \ match_1$\mbox{\tt \ ;}\\
   &\mbox{\tt \ \ \ \ \ \ \ \ \ \ \ \ \ \ \ \ {\char'137}\ \ ->\ }$\ldots{}$\mbox{\tt \ case\ }$v$\mbox{\tt \ of\ {\char'173}}\\
   &\mbox{\tt \ \ \ \ \ \ \ \ \ \ \ \ \ \ \ \ \ \ \ \ \ \ \ \ \ \ \ }$p_n\ \ match_n$\ \mbox{\tt ;}\\
   &\mbox{\tt \ \ \ \ \ \ \ \ \ \ \ \ \ \ \ \ \ \ \ \ \ \ \ \ \ \ \ {\char'137}\ \ ->\ error\ "No\ match"\ {\char'175}}$\ldots$\mbox{\tt {\char'175}}\\
   &\mbox{\tt \ }{\rm where each $match_i$ has the form:}\\
   &\mbox{\tt \ \ |\ }$gs_{i,1}$  \mbox{\tt \ ->\ }$e_{i,1}$\mbox{\tt \ ;\ }$\ldots{}$\mbox{\tt \ ;\ |\ }$gs_{i,m_i}$\mbox{\tt \ ->\ }$e_{i,m_i}$\mbox{\tt \ where\ {\char'173}\ }$decls_i$\mbox{\tt \ {\char'175}}\\[4pt]
%\\
(c)&\mbox{\tt case\ }$v$\mbox{\tt \ of\ {\char'173}\ }$p$\mbox{\tt \ |\ }$gs_1$\mbox{\tt \ ->\ }$e_1$\mbox{\tt \ ;\ }$\ldots{}$\\
   &\hspace*{4pt}\mbox{\tt \ \ \ \ \ \ \ \ \ \ \ \ \ |\ }$gs_n$\mbox{\tt \ ->\ }$e_n$\mbox{\tt \ where\ {\char'173}\ }$decls$\mbox{\tt \ {\char'175}}\\
   &\hspace*{2pt}\mbox{\tt \ \ \ \ \ \ \ \ \ \ \ \ {\char'137}\ \ \ \ \ ->\ }$e'$\mbox{\tt \ {\char'175}}\\
   &$=$\mbox{\tt \ case\ }$e'$\mbox{\tt \ of\ {\char'173}\ }$y$\mbox{\tt \ ->}\\
   &\mbox{\tt \ \ \ case\ }$v$\mbox{\tt \ of\ {\char'173}}\\
   &\mbox{\tt \ \ \ \ \ }$p$\mbox{\tt \ ->\ let\ {\char'173}\ }$decls$\mbox{\tt \ {\char'175}\ in}\\
   &\hprime{\mbox{\tt \ \ \ \ \ \ \ \ \ \ case\ ()\ of\ {\char'173}}}\\
   &\hprime{\mbox{\tt \ \ \ \ \ \ \ \ \ \ \ \ ()\ |\ }$gs_1$\mbox{\tt \ ->\ }$e_1$\mbox{\tt ;}}\\
   &\hprime{\mbox{\tt \ \ \ \ \ \ \ \ \ \ \ \ {\char'137}\ ->\ }$\ldots$\mbox{\tt \ case\ ()\ of\ {\char'173}}}\\
   &\hprime{\mbox{\tt \ \ \ \ \ \ \ \ \ \ \ \ \ \ \ \ \ \ \ \ \ \ \ ()\ |\ }$gs_n$\mbox{\tt \ ->\ }$e_n$\mbox{\tt ;}}\\
   &\hprime{\mbox{\tt \ \ \ \ \ \ \ \ \ \ \ \ \ \ \ \ \ \ \ \ \ \ \ {\char'137}\ \ ->\ }$y$\mbox{\tt \ {\char'175}\ }$\ldots$\mbox{\tt \ {\char'175}}}\\
   &\mbox{\tt \ \ \ \ \ {\char'137}\ ->\ }$y$\mbox{\tt \ {\char'175}{\char'175}}\\
   &{\rm where $y$ is a new variable}\\[4pt]
%\\
(d)&\mbox{\tt case\ }$v$\mbox{\tt \ of\ {\char'173}\ {\char'176}}$p$\mbox{\tt \ ->\ }$e$\mbox{\tt ;\ {\char'137}\ ->\ }$e'$\mbox{\tt \ {\char'175}}\\
&$=$\mbox{\tt \ ({\char'134}}$x_1$ $\ldots$ $x_n$ \mbox{\tt ->} $e$ \mbox{\tt )\ (case\ }$v$\mbox{\tt \ of\ {\char'173}\ }$p$\mbox{\tt ->} 
$x_1$\mbox{\tt \ {\char'175})} $\ldots$ \mbox{\tt (case\ }$v$\mbox{\tt \ of\ {\char'173}\ }$p$\mbox{\tt \ ->\ }$x_n$\mbox{\tt {\char'175})}\\
&{\rm where $x_1, \ldots, x_n$ are all the variables in $p\/$}\\[4pt]
%\\
(e)&\mbox{\tt case\ }$v$\mbox{\tt \ of\ {\char'173}\ }$x${\tt @}$p$\mbox{\tt \ ->\ }$e$\mbox{\tt ;\ {\char'137}\ ->\ }$e'$\mbox{\tt \ {\char'175}}\\
&$=$\mbox{\tt \ \ case\ }$v$\mbox{\tt \ of\ {\char'173}\ }$p$\mbox{\tt \ ->\ (\ {\char'134}\ }$x$\mbox{\tt \ ->\ }$e$\mbox{\tt \ )\ }$v$\mbox{\tt \ ;\ {\char'137}\ ->\ }$e'$\mbox{\tt \ {\char'175}}\\[4pt]
%\\
(f)&\mbox{\tt case\ }$v$\mbox{\tt \ of\ {\char'173}\ {\char'137}\ ->\ }$e$\mbox{\tt ;\ {\char'137}\ ->\ }$e'$\mbox{\tt \ {\char'175}\ }$=$\mbox{\tt \ }$e$\\[4pt]
\end{tabular}
}
%**<div align=center> <h4>Figure 3.1</h4> </div>
\ecaption{Semantics of Case Expressions, Part 1}
\label{simple-case-expr-1}
\end{figure}

\begin{figure}[tb]
\outlinec{\small
\begin{tabular}{@{}cl}
(g)&\mbox{\tt case\ }$v$\mbox{\tt \ of\ {\char'173}\ }$K\ p_1 \ldots p_n$\mbox{\tt \ ->\ }$e$\mbox{\tt ;\ {\char'137}\ ->\ }$e'$\mbox{\tt \ {\char'175}}\\
&$=$\mbox{\tt \ case\ }$v$\mbox{\tt \ of\ {\char'173}}\\
&\mbox{\tt \ \ \ \ \ }$K\ x_1 \ldots x_n$\mbox{\tt \ ->\ case\ }$x_1$\mbox{\tt \ of\ {\char'173}}\\
&\mbox{\tt \ \ \ \ \ \ \ \ \ \ \ \ \ \ \ \ \ \ \ \ }$p_1$\mbox{\tt \ ->\ }$\ldots$\mbox{\tt \ case\ }$x_n$\mbox{\tt \ of\ {\char'173}\ }$p_n$\mbox{\tt \ ->\ }$e$\mbox{\tt \ ;\ {\char'137}\ ->\ }$e'$\mbox{\tt \ {\char'175}\ }$\ldots$\\
&\mbox{\tt \ \ \ \ \ \ \ \ \ \ \ \ \ \ \ \ \ \ \ \ {\char'137}\ \ ->\ }$e'$\mbox{\tt \ {\char'175}}\\
&\mbox{\tt \ \ \ \ \ {\char'137}\ ->\ }$e'$\mbox{\tt \ {\char'175}}\\[2pt]
&{\rm at least one of $p_1, \ldots, p_n$ is not a variable; $x_1, \ldots, x_n$ are new variables}\\[4pt]
%\\
(h)&\mbox{\tt case\ }$v$\mbox{\tt \ of\ {\char'173}\ }$k$\mbox{\tt \ ->\ }$e$\mbox{\tt ;\ {\char'137}\ ->\ }$e'$\mbox{\tt \ {\char'175}\ }$=$\mbox{\tt \ if\ (}$v$\mbox{\tt ==}$k$\mbox{\tt )\ then\ }$e$\mbox{\tt \ else\ }$e'$ \\
&{\rm where $k$ is a numeric, character, or string literal} \\[4pt]
%\\
(i)&\mbox{\tt case\ }$v$\mbox{\tt \ of\ {\char'173}\ }$x$\mbox{\tt \ ->\ }$e$\mbox{\tt ;\ {\char'137}\ ->\ }$e'$\mbox{\tt \ {\char'175}\ }$=$\mbox{\tt \ case\ }$v$\mbox{\tt \ of\ {\char'173}\ }$x$\mbox{\tt \ ->\ }$e$\mbox{\tt \ {\char'175}}\\[4pt]
%\\
(j)&\mbox{\tt case\ }$v$\mbox{\tt \ of\ {\char'173}\ }$x$\mbox{\tt \ ->\ }$e$\mbox{\tt \ {\char'175}\ }$=$\mbox{\tt \ (\ {\char'134}\ }$x$\mbox{\tt \ ->\ }$e$\mbox{\tt \ )\ }$v$\\[4pt]
%\\
(k)&\mbox{\tt case\ }$N$ $v$\mbox{\tt \ of\ {\char'173}\ }$N$\mbox{\tt \ }$p$\mbox{\tt \ ->\ }$e$\mbox{\tt ;\ {\char'137}\ ->\ }$e'$\mbox{\tt \ {\char'175}}\\
&$=$\mbox{\tt \ case\ }$v$\mbox{\tt \ of\ {\char'173}\ }$p$\mbox{\tt \ ->\ }$e$\mbox{\tt ;\ {\char'137}\ ->\ }$e'$\mbox{\tt \ {\char'175}}\\
&{\rm where $N$ is a \mbox{\tt newtype} constructor}\\[4pt]
(l)&\mbox{\tt case\ }$\bot$\mbox{\tt \ of\ {\char'173}\ }$N$\mbox{\tt \ }$p$\mbox{\tt \ ->\ }$e$\mbox{\tt ;\ {\char'137}\ ->\ }$e'$\mbox{\tt \ {\char'175}\ }$=$\mbox{\tt \ case\ }$\bot$\mbox{\tt \ of\ {\char'173}\ }$p$\mbox{\tt \ ->\ }$e$\mbox{\tt \ {\char'175}}\\
&{\rm where $N$ is a \mbox{\tt newtype} constructor}\\[4pt]
(m)& \mbox{\tt case\ } $v$ \mbox{\tt \ of\ {\char'173}\ } $K$ \mbox{\tt \ {\char'173}} $f_1$ \mbox{\tt \ =\ } $p_1$ \mbox{\tt \ ,\ } $f_2$ \mbox{\tt \ =\ }
$p_2$ \mbox{\tt \ ,\ } $\ldots$ \mbox{\tt {\char'175}\ ->\ } $e$ \mbox{\tt ;\ {\char'137}\ ->\ } $e'$ \mbox{\tt \ {\char'175}}\\
&$=$ \mbox{\tt \ case\ }$e'$\mbox{\tt \ of\ {\char'173}}\\
&\mbox{\tt \ \ \ }$y$\mbox{\tt \ ->}\\
&\mbox{\tt \ \ \ \ case\ } $v$ \mbox{\tt \ of\ {\char'173}\ }\\
&\mbox{\tt \ \ \ \ \ } $K$ \mbox{\tt \ {\char'173}\ } $f_1$ \mbox{\tt \ =\ } $p_1$ \mbox{\tt \ {\char'175}\ ->}\\
&\mbox{\tt \ \ \ \ \ \ \ \ \ \ \ \ case\ } $v$ \mbox{\tt \ of\ {\char'173}} $K$ \mbox{\tt \ {\char'173}} $f_2$ \mbox{\tt \ =\ } $p_2$ \mbox{\tt \ ,\ }
$\ldots$ \mbox{\tt \ {\char'175}\ ->\ } $e$ \mbox{\tt ;\ {\char'137}\ ->\ } $y$ \mbox{\tt \ {\char'175};}\\
&\mbox{\tt \ \ \ \ \ \ \ \ \ \ \ \ {\char'137}\ ->\ } $y$ \mbox{\tt \ {\char'175}{\char'175}}\\
&{\rm where $f_1$, $f_2$, $\ldots$ are fields of constructor $K$; $y$
is a new variable}\\[4pt]
(n)&\mbox{\tt case\ } $v$ \mbox{\tt \ of\ {\char'173}\ } $K$ \mbox{\tt \ {\char'173}} $f$ \mbox{\tt \ =\ } $p$ \mbox{\tt {\char'175}\ ->\ } $e$ \mbox{\tt ;\ {\char'137}\ ->\ }
$e'$ \mbox{\tt \ {\char'175}} \\
&$=$\mbox{\tt \ case\ } $v$ \mbox{\tt \ of\ {\char'173}}\\
&   \mbox{\tt \ \ \ \ } $K$ $p_1\ \ldots \  p_n$ \mbox{\tt \ ->\ } $e$ \mbox{\tt ;\ {\char'137}\ ->\ } $e'$ \mbox{\tt \ {\char'175}}\\
&{\rm where $p_i$ is $p$ if $f$ labels the $i$th component of $K$,
\mbox{\tt {\char'137}} otherwise}\\
(o)&\mbox{\tt case\ } $v$ \mbox{\tt \ of\ {\char'173}\ } $K$ \mbox{\tt \ {\char'173}{\char'175}\ ->\ } $e$ \mbox{\tt ;\ {\char'137}\ ->\ }
$e'$ \mbox{\tt \ {\char'175}} \\
&$=$\mbox{\tt \ case\ } $v$ \mbox{\tt \ of\ {\char'173}}\\
&   \mbox{\tt \ \ \ \ } $K$ \mbox{\tt {\char'137}} $\ldots$ \mbox{\tt {\char'137}\ ->\ } $e$ \mbox{\tt ;\ {\char'137}\ ->\ } $e'$ \mbox{\tt \ {\char'175}}\\
(p)&\mbox{\tt case\ (}$K'$\mbox{\tt \ }$e_1$\mbox{\tt \ }$\ldots$\mbox{\tt \ }$e_m$\mbox{\tt )\ of\ {\char'173}\ }$K$\mbox{\tt \ }$x_1$\mbox{\tt \ }$\ldots$\mbox{\tt \ }$x_n$\mbox{\tt \ ->\ }$e$\mbox{\tt ;\ {\char'137}\ ->\ }$e'$\mbox{\tt \ {\char'175}\ }$=$\mbox{\tt \ }$e'$\\
&{\rm where $K$ and $K'$ are distinct \mbox{\tt data} constructors of arity $n$ and $m$, respectively}\\[4pt]
%\\
(q)&\mbox{\tt case\ (}$K$\mbox{\tt \ }$e_1$\mbox{\tt \ }$\ldots$\mbox{\tt \ }$e_n$\mbox{\tt )\ of\ {\char'173}\ }$K$\mbox{\tt \ }$x_1$\mbox{\tt \ }$\ldots$\mbox{\tt \ }$x_n$\mbox{\tt \ ->\ }$e$\mbox{\tt ;\ {\char'137}\ ->\ }$e'$\mbox{\tt \ {\char'175}}\\
&$=$\mbox{\tt \ ({\char'134}}$x_1~\ldots~x_n$\mbox{\tt \ ->\ }$e$\mbox{\tt )\ }$e_1~\ldots~e_n$\\
&{\rm where $K$ is a \mbox{\tt data} constructor of arity $n$}\\[4pt]

(r)&\mbox{\tt case}~$\bot$~\mbox{\tt of\ {\char'173}\ }$K$\mbox{\tt \ }$x_1$\mbox{\tt \ }$\ldots$\mbox{\tt \ }$x_n$\mbox{\tt \ ->\ }$e$\mbox{\tt ;\ {\char'137}\ ->\ }$e'$\mbox{\tt \ {\char'175}} ~$=$~ $\bot$ \\
&{\rm where $K$ is a \mbox{\tt data} constructor of arity $n$}\\[4pt]
\end{tabular}
}
%**<div align=center> <h4>Figure 3.2</h4> </div>
\ecaption{Semantics of Case Expressions, Part 2}
\label{simple-case-expr-2}
\end{figure}

\begin{figure}[tb]
\outlinec{\small
\begin{haskellprime}
\begin{tabular}{@{}cl}
(s)&\mbox{\tt case\ ()\ of\ {\char'173}\ ()\ |\ }$g_1$\mbox{\tt ,\ }$\ldots$\mbox{\tt ,\ }$g_n$\mbox{\tt \ ->\ }$e$\mbox{\tt ;\ {\char'137}\ ->\ }$e'$\mbox{\tt \ {\char'175}}\\
   &$=$\mbox{\tt \ case\ ()\ of\ {\char'173}}\\
   &\mbox{\tt \ \ \ \ ()\ |\ }$g_1$\mbox{\tt \ ->\ }\ldots\mbox{\tt \ case\ ()\ of\ {\char'173}}\\
   &\mbox{\tt \ \ \ \ \ \ \ \ \ \ \ \ \ \ \ \ \ \ \ ()\ |\ }$g_n$\mbox{\tt \ ->\ }$e$\mbox{\tt ;}\\
   &\mbox{\tt \ \ \ \ \ \ \ \ \ \ \ \ \ \ \ \ \ \ \ {\char'137}\ ->\ }$e'$\mbox{\tt \ {\char'175}\ }\ldots\\
   &\mbox{\tt \ \ \ \ {\char'137}\ ->\ }$e'$\mbox{\tt \ {\char'175}}\\
   &{\rm where $y$ is a new variable}\\[4pt]

(t)&\mbox{\tt case\ ()\ of\ {\char'173}\ ()\ |\ }$p$\mbox{\tt \ <-\ }$e_0$\mbox{\tt \ ->\ }$e$\mbox{\tt ;\ {\char'137}\ ->\ }$e'$\mbox{\tt \ {\char'175}}\\
   &$=$\mbox{\tt \ case\ }$e_0$\mbox{\tt \ of\ {\char'173}\ }$p$\mbox{\tt \ ->\ }$e$\mbox{\tt ;\ {\char'137}\ ->\ }$e'$\mbox{\tt \ {\char'175}}\\

(u)&\mbox{\tt case\ ()\ of\ {\char'173}\ ()\ |\ let\ }$decls$\mbox{\tt \ ->\ }$e$\mbox{\tt ;\ {\char'137}\ ->\ }$e'$\mbox{\tt \ {\char'175}}\\
   &$=$\mbox{\tt \ let\ }$decls$\mbox{\tt \ in\ }$e$\\[4pt]

(v)&\mbox{\tt case\ ()\ of\ {\char'173}\ ()\ |\ }$e_0$\mbox{\tt \ ->\ }$e$\mbox{\tt ;\ {\char'137}\ ->\ }$e'$\mbox{\tt \ {\char'175}}\\
   &$=$\mbox{\tt \ if\ }$e_0$\mbox{\tt \ then\ }$e$\mbox{\tt \ else\ }$e'$\\[4pt]
\end{tabular}
\end{haskellprime}
}
%**<div align=center> <h4>Figure 3.3</h4> </div>
\ecaption{\hprime{Semantics of Case Expressions, Part 3}}
\label{simple-case-expr-3}
\end{figure}




In Figures~\ref{simple-case-expr-1}--\ref{simple-case-expr-3}:
\mbox{$\it e$}, \mbox{$\it e'$} and \mbox{$\it e_i$} are expressions; 
\mbox{$\it g_i$} and \mbox{$\it gs_i$} are guards and sequences of guards respecively;
\mbox{$\it p$} and \mbox{$\it p_i$} are patterns; 
\mbox{$\it v$}, \mbox{$\it x$}, and \mbox{$\it x_i$} are variables; 
\mbox{$\it K$} and \mbox{$\it K'$} are algebraic datatype (\mbox{\tt data}) constructors (including
tuple constructors);  and \mbox{$\it N$} is a \mbox{\tt newtype} constructor\index{newtype declaration@{\tt newtype} declaration}.
% For clarity, several rules are expressed using
% \mbox{\tt let} (used only in a non-recursive way); their usual purpose is to
% prevent name capture
% (e.g., in rule~(c)).  The rules may be re-expressed entirely with
% \mbox{\tt cases} by applying this identity:
% \[
% \mbox{$\it \makebox{\tt let\ }x\makebox{\tt \ =\ }y\makebox{\tt \ in\ }e\makebox{\tt \ }\ =\makebox{\tt \ \ case\ }y\makebox{\tt \ of\ {\char'173}\ }x\makebox{\tt \ ->\ }e\makebox{\tt \ {\char'175}}$}
% \]

%Using all but the last two identities (rules~(n) and~(o)) in Figure~\ref{simple-case-expr-2}
%in a left-to-right manner yields a translation into a
%subset of general \mbox{\tt case} expressions called {\em simple case expressions}.%
%\index{simple case expression}
Rule~(b) matches a general source-language
\mbox{\tt case} expression, regardless of whether it actually includes
guards---if no guards are written, then \mbox{\tt True} is substituted for the guards \mbox{$\it gs_{i,j}$}
in the \mbox{$\it match_i$} forms.
Subsequent identities manipulate the resulting \mbox{\tt case} expression into simpler
and simpler forms.
%The semantics of simple \mbox{\tt case} expressions is 
%given by the last two identities ((n) and~(o)).

Rule~(h) in Figure~\ref{simple-case-expr-2} involves the
overloaded operator \mbox{\tt ==}; it is this rule that defines the
meaning of pattern matching against overloaded constants.
\index{pattern-matching!overloaded constant}

%% When used as a translation, the identities in
%% Figure~\ref{simple-case-expr} will generate a very inefficient
%% program.  This can be fixed by using further \mbox{\tt case} or \mbox{\tt let} 
%% expressions, but doing so 
%% would clutter the identities, which are intended only to convey the semantics.

These identities all preserve the static semantics.  Rules~(d), (e), (j), and~(q)
use a lambda rather than a \mbox{\tt let}; this indicates that variables bound
by \mbox{\tt case} are monomorphically typed (Section~\ref{type-semantics}).
\index{monomorphic type variable}

%**~footer

% Local Variables: 
% mode: latex
% End:
